%% 第26回編集委員会（物理第4巻）：組 電W1
%% p4-02.tex  電-2　電気力線を数える：ガウスの法則と電位（★3〜5・3段）
%% 中身は計画書 tools/phys4-plan.md の B「電-2」。公式の持ち主は A（A22〜A25・A61・A62）。
%% 3段：目でわかる（①②、★3）→ 手でわかる（③〜⑦、★3。ここまでで十分）→ 式でわかる（⑧⑨ ★4、⑩〜⑫ ★5）。
%%   X-1・Y-1 のとおり、ガウスの法則を「袋の積分（∯）」で書く所と grad・div・rot は★4、境界値問題・一意性・電界解析は★5。
%%   手でわかるの段は、高校の「電気力線の本数」と「グラフの面積（仕事）」だけでたどれるように書いた。
%% 数値はすべて Python で計算した：checks/物理第4巻/電-2/calc_p4_02.py（出力 calc_p4_02_out.txt）、
%%   fig_p4_02.py（⑦の等電位線と電気力線の座標、⑫のとがった電極のラプラス方程式と等電位線）。
%% 数学との約束（計画書 C-6-2）：
%%   ・8-8:1010「この本と物理シリーズでは ∯E·n dS を電場の流束 Φ_E とよぶ」→ ③⑧は Φ_E で書き、電束密度 D・電束は専門書の囲みで読み方だけ（D は電-3 ⑧）。
%%   ・8-11⑦（ガウスの法則は袋を選ぶ。球・線・板）→ ④は本数の数え方で同じ答えを出し、計算は 8-11⑦ を指す（Y-9）。
%%   ・8-11:483（板の σ/(2ε0)）→ ④。8-12:549 の点電荷の電位 φ＝Q/(4πε0)×1/r（「物理シリーズ『電磁気』」は『電磁気学』の誤り：報告に書く）→ ⑤。
%%   ・8-12③（板の温度。42.857・52.679・25）→ ⑪・例題11（同じ数）。8-4⑧（E＝−grad φ、0.01 m・100 V・10000 V/m）→ ⑥⑨・例題4。
%%   ・8-9⑬（div E＝ρ/ε0）→ ⑧。8-7⑨（保存力・∮）→ ⑤⑨。
%% 前の巻との約束：計画書 A-6 に電-2 の行はない。波-7 ②の逆2乗の法則（点の音源）を、はてな？の答えで「同じ理由」として指した。
%%   力-4 の保存力と位置エネルギー（F＝−dU/dx）を⑤⑨で使った。
%% 記号（この章）：電位 φ（数学 8-4 と同じ。高校の電位 V）、電圧（電位差）V、電場の流束 Φ_E、電荷密度 ρ（この章では抵抗率は出ない）、
%%   位置エネルギー U、仕事 W、袋の中の電荷 Q_中、面積 S（Y-4）、球の半径 R、同軸の内・外の半径 a・b、板の間かく d、電子ボルト eV。
%%   立体角は記号を作らない（数学 8-8 の Ω を一言で指すだけ）。
%% 出典（11月5日に確かめた）
%%   ・CODATA 2022（電-1 と同じ。ε0、e、mₑ、mₚ）。
%%   ・晴天の大気電場 約100 V/m（地面の近く、下向き）、電離層と地面の電位の差 約300 kV、雷雲1つが運ぶ電荷 約200 C（全地球の回路の見積もり）：
%%       D. Siingh ほか "The atmospheric global electric circuit: An overview", Atmos. Res. 84, 91 (2007)、arXiv:0906.1280（本文を pdftotext で読んだ。
%%       「fair weather electric field is ≈ 10² V/m at the ground level」「~300 kV between the ground and the ionosphere」）。
%%       National Academies Press "The Earth's Electrical Environment"（1986）の晴天の電場 100〜150 V/m（WebSearch の抜粋）と照らし合わせた。
%%   ・空気がこわれる（放電が始まる）電場 約3×10⁶ V/m（海面の高さ）、雲の中で測られる電場はそれよりずっと小さい（約10万〜15万 V/m）：
%%       J. R. Dwyer への取材記事（Physics World）と、雷の教科書の章（Springer、DOI 10.1007/978-94-017-8938-7_7）の WebSearch の抜粋で照らし合わせた。
%%   ・雷の始まりのしくみがわかっていないこと：J. R. Dwyer, M. A. Uman, "The physics of lightning", Phys. Rep. 534, 147–241 (2014)、
%%       DOI 10.1016/j.physrep.2013.09.004（要旨の「雷の雲の中で何が雷を始めるのか」が未解決の問題に挙がっていることを WebSearch の抜粋で確かめた。本文は未読）。
%%       宇宙線がつくる高速の電子と、乱れた雲の中の小さな強い電場の所を組み合わせた説：A. Yu. Kostinskiy, T. C. Marshall, M. Stolzenburg,
%%       J. Geophys. Res. Atmos. (2020)、DOI 10.1029/2020JD033191、arXiv:1906.01033（要旨を arxiv.org で読んだ。「a possible mechanism」として提案）。
%% 〔要確認〕なし。
\makeatletter
% 点電荷（丸の中に＋・−。色だけで分けない）  \sxcplus{(中心)}{半径}  \sxcminus{(中心)}{半径}
\newcommand\sxcplus[2]{\filldraw[fill=bkRed!75!black, draw=bkInk, line width=0.5pt] #1 circle (#2); \node[font=\sffamily\bfseries\footnotesize, text=white, inner sep=0pt] at #1 {$+$};}
\newcommand\sxcminus[2]{\filldraw[fill=bkBlue!75!black, draw=bkInk, line width=0.5pt] #1 circle (#2); \node[font=\sffamily\bfseries\footnotesize, text=white, inner sep=0pt] at #1 {$-$};}
\newcommand\sxcp[1]{\node[font=\sffamily\bfseries\footnotesize, text=bkRed!75!black, inner sep=0pt] at #1 {$+$};}
\newcommand\sxcm[1]{\node[font=\sffamily\bfseries\footnotesize, text=bkBlue!75!black, inner sep=0pt] at #1 {$-$};}
\tikzset{sxc lab/.style={font=\sffamily\footnotesize, text=bkInk},
  sxc fl/.style={bkPurple!80!black, line width=0.7pt},
  sxc tip/.style={bkPurple!80!black, line width=0.7pt, -{Stealth[length=1.9mm, width=1.5mm]}},
  sxc eq/.style={bkGreen!50!black, line width=0.8pt, dash pattern=on 2.4pt off 1.4pt},
  sxc lap/.style={bkGreen!50!black, line width=0.6pt},
  sxc arr/.style={bkNavy, line width=0.8pt, -{Stealth[length=1.8mm, width=1.4mm]}},
  sxc f/.style={line width=1.2pt, -{Stealth[length=2.6mm, width=2.0mm]}},
  sxc dim/.style={bkGray, line width=0.5pt, {Stealth[length=1.6mm]}-{Stealth[length=1.6mm]}}}
\makeatother

\phchapter{2}{電気力線を数える：\\ガウスの法則と電位}{電気の力は、なぜはなれるほど弱くなる？}{\bklevelrange{3}{5}}{電気力線を数える：ガウスの法則と電位}
\bkmichishirube{4-5,5-10,8-4,8-7,8-8,8-9,8-11,8-12,力-4,電-1}{電-3,電-4,電-6,電-8,電-14}

\begin{goal}
\begin{itemize}[label={\color{bkTheme}●}, leftmargin=1.6\zw, itemsep=0pt, topsep=1pt]
\item \bklevel{3}\enspace ガウスの法則で、球・長い直線・広い平面の電荷の電場を求められる。
\item \bklevel{3}\enspace 電位を使って、電荷を運ぶ仕事や、加速された電子の速さを求められる。
\item \bklevel{3}\enspace 等電位面と電気力線の関係を読み、一様な電場で $V=Ed$ を使える。
\item \bklevel{4}\enspace ガウスの法則の微分形と $\vec E=-\mathrm{grad}\,\phi$ を、ベクトル解析で説明できる。
\item \bklevel{5}\enspace ラプラス方程式を表計算で解き、電場が集まる所を示せる。
\end{itemize}
\end{goal}

\phkoushiki{%
\phk{電荷 $Q$ から出る電気力線の本数 $4\pi kQ$（$=Q/\varepsilon_0$）}{高校物理}{ph:電-2:honsu}
\phk{点電荷の電位 $\phi=k\,\dfrac{Q}{r}$（教科書の $V$）、位置エネルギー $U=q\phi$、仕事 $W=qV$}{高校物理}{ph:電-2:deni}
\phk{一様な電場 $V=Ed$}{高校物理}{ph:電-2:ved}
\phk{等電位面と電気力線は垂直に交わる}{高校物理}{ph:電-2:suichoku}
}

\begin{hatena}
\noindent\begin{minipage}[c]{0.55\linewidth}\raggedright
こすった下じきを髪の毛に近づけると、髪は立ち上がりました（電-1）。でも、下じきを遠ざけると、髪はすぐに下がります。電気の力は、離れるほど弱くなるのです。電-1 では、それを「距離の2乗に反比例する」というクーロンの法則で表しました。

では、なぜ2乗なのでしょう。3乗でも1乗でもなく、ちょうど2乗になるのには、わけがあるのでしょうか。
\end{minipage}\hfill
\begin{minipage}[c]{0.42\linewidth}\centering
\begin{tikzpicture}[x=0.85mm, y=0.85mm]
  \foreach \a in {0,30,...,330} {\draw[sxc fl] (\a:2.4) -- (\a:25);}
  \draw[bkOrange!80!black, line width=0.8pt] (0,0) circle (10);
  \draw[bkOrange!80!black, line width=0.8pt, dash pattern=on 2.4pt off 1.4pt] (0,0) circle (20);
  \sxcplus{(0,0)}{2.4}
  \node[sxc lab, fill=bkYellow!12, inner sep=1pt] at (0,-12.8) {$r$};
  \node[sxc lab, fill=bkYellow!12, inner sep=1pt] at (0,-22.8) {$2r$};
  \node[sxc lab, anchor=south, fill=bkYellow!12, inner sep=1pt] at (0,25.5) {どちらの円も、線は同じ数};
\end{tikzpicture}
\end{minipage}
\end{hatena}

\noindent 電-1 では、電荷の周りの空間に\textbf{電場}があり、電場の向きをつないだ線が\textbf{電気力線}でした。この章では、電気力線の「本数」を数えるという、とても見通しのよい考え方を学びます。本数のきまりは\textbf{ガウスの法則}とよばれ、クーロンの法則と同じ内容を、別の形で言い表したものです。後半では、電気の力がする仕事から\textbf{電位}（電気の高さ）を作ります。電池の「1.5\,V」や、家のコンセントの「100\,V」の V は、この電位の差の単位です。

\begin{yattemiyou}[光の広がりと、電場の向き]
\textbf{用意するもの}：スマホ（ライトと、明るさをはかるアプリ）、ものさし、アルミはく、糸、ストロー、ティッシュ\par\vspace{2pt}
\noindent\begin{minipage}[c]{0.64\linewidth}
\begin{enumerate}[label={\bkcircled[bkGreen]{\arabic*}}, leftmargin=2.4\zw, itemsep=1pt, topsep=0pt]
\item 暗い部屋で、スマホのライトを壁に向け、壁から\qty{20}{cm}と\qty{40}{cm}で、壁にできる光の輪の直径をはかる。
\item 別のスマホのアプリで、ライトから\qty{20}{cm}と\qty{40}{cm}の所の明るさをはかる。
\item アルミはくを丸めた玉を糸でつるし、ティッシュでこすったストローを触れさせる。細かくちぎったティッシュを、玉の周りのいろいろな所に近づける。
\end{enumerate}
\end{minipage}\hfill
\begin{minipage}[c]{0.33\linewidth}\centering
\begin{tikzpicture}[x=0.8mm, y=0.8mm]
  \fill[bkYellow!30] (0,10) -- (36,0) -- (36,20) -- cycle;
  \draw[bkInk, line width=0.4pt] (0,10) -- (36,0) (0,10) -- (36,20);
  \draw[bkGray, line width=0.4pt, dashed] (18,5) -- (18,15);
  \fill[bkInk] (-2,8) rectangle (1,12);
  \draw[bkInk, line width=1pt] (36,-1) -- (36,21);
  \draw[sxc dim] (0,-3) -- (18,-3);
  \node[sxc lab, anchor=north] at (9,-3.5) {\qty{20}{cm}};
  \draw[sxc dim] (0,-9) -- (36,-9);
  \node[sxc lab, anchor=north] at (18,-9.5) {\qty{40}{cm}};
\end{tikzpicture}
\end{minipage}
\end{yattemiyou}

\noindent やってみようの答え：①\qty{40}{cm}では、光の輪の直径が\qty{20}{cm}のときの約2倍、面積は約4倍になります。②明るさは約 $\tfrac14$ になります（ライトの光が円すいの形に広がり、同じ光が4倍の面積に分かれるから）。③ティッシュの切れはしは、どこから近づけても、玉の中心へ向かって飛びつきます。玉の周りの電場は、玉の中心から放射状に（またはその向きへ）向いているからです。①②の「同じ量が広がると、面積の分だけ薄まる」という考えが、この章のはてな？の答えのかぎになります。

\begin{omoidasou}
\bkrecall{電-1}{点電荷の電場 $E=kQ/r^2$、$k=1/(4\pi\varepsilon_0)$。電気力線は正の電荷から出て負の電荷に入り、こみ合っている所ほど電場が強い。}
\bkrecall{4-5}{半径 $r$ の球の表面積は $4\pi r^2$、円柱の側面の面積は $2\pi r\times$（長さ）。}
\bkrecall{力-4}{仕事＝力×力の向きに動いた距離。保存力の仕事は道によらず、位置エネルギー $U$ の差で決まる。力は $F=-\dd U/\dd x$。}
\bkrecall{5-10}{グラフの面積は積分で求める。$\displaystyle\int_r^{\infty}\frac{\dd x}{x^2}=\frac{1}{r}$。}
\bkrecall{8-8}{面を通りぬける量を流束という。点電荷を囲む袋の電場の流束は、袋の形によらず $Q/\varepsilon_0$（8-8⑫）。}
\bkrecall{8-12}{電荷のない所の電位は、どの点も周りの平均と同じ（ラプラス方程式。8-12③の板の温度と同じ）。}
\end{omoidasou}

\noindent この章の記号：\textbf{電位}を $\phi$（ファイ）と書きます（数学 8-4・8-12 と同じ）。高校の教科書は電位を $V$ と書きますが、この本では、2点の電位の差（\textbf{電圧}）を $V$ と書き、電位そのものとは文字で分けます。電荷の密度（\qty{1}{m^3}あたりの電荷）は $\rho$（ロー）と書きます。

%% ======================== 物理の道具 ========================
\bksection{doguu}{物理の道具}

\phstage{me}{電気力線を数える}{3}

\bksubsection{1}{電気力線の本数を数える}
\label{ph:電-2:sec1}
電-1 では、電気力線は「こみ合っている所ほど電場が強い」と決めました。ここでは、それを数で約束します。
\begin{bkcenter}
\fcolorbox{bkTheme}{white}{\parbox{0.86\linewidth}{\centering 電場に垂直な面 \qty{1}{m^2} を通りぬける電気力線の本数を、電場の強さ $E$〔N/C〕に等しくかく}}
\end{bkcenter}
電場が\qty{899}{N/C}の所では、電場に垂直な面 \qty{1}{m^2} あたり899本の線を通す、と考えるのです。もちろん、本物の線があるわけではありません。電場の強さを、線の「こみぐあい」で見えるようにするための約束です。

\phboxhead{点電荷から出る本数}
この約束で、電荷 $Q$ の点電荷から、全部で何本の線が出るかを数えましょう。電荷を中心とする半径 $r$ の球面を考えます。球面の上では、電場はどこでも球面に垂直（外向き）で、強さはどこでも $E=kQ/r^2$ です。球面の面積は $4\pi r^2$ なので、球面を通りぬける本数は
\pheq{N=E\times4\pi r^2=k\,\frac{Q}{r^2}\times4\pi r^2=4\pi kQ=\frac{Q}{\varepsilon_0}}{N·m²/C ＝ N/C × m²}
\label{ph:電-2:honsu}です（$k=1/(4\pi\varepsilon_0)$ を使った）。\textbf{$r$ が消えました}。半径 \qty{0.1}{m}の球面でも、\qty{1}{km}の球面でも、通る本数は同じ $Q/\varepsilon_0$ 本です。電荷から出た線は、とちゅうで増えも減りもせず、遠くまで届くということです。\qty{1}{nC}の電荷なら $10^{-9}\div(8.85\times10^{-12})\approx113$ 本です。高校の教科書の「電荷 $Q$ から出る電気力線は $4\pi kQ$ 本」は、このことです。\bkidx{でんきりきせんのほんすう}{電気力線の本数}

\begin{bkfigure}
\begin{tikzpicture}[x=1mm, y=1mm]
  \foreach \a in {0,22.5,...,337.5} {\draw[sxc fl] (\a:2.4) -- (\a:29);}
  \draw[bkOrange!80!black, line width=0.9pt] (0,0) circle (12);
  \draw[bkOrange!80!black, line width=0.9pt, dash pattern=on 2.4pt off 1.4pt] (0,0) circle (24);
  \sxcplus{(0,0)}{2.4}
  \draw[bkRed!80!black, line width=2.4pt] (-10:12) arc (-10:10:12);
  \draw[bkRed!80!black, line width=2.4pt] (-10:24) arc (-10:10:24);
  \node[sxc lab, fill=white, inner sep=1pt] at (-101:12) {$r$};
  \node[sxc lab, fill=white, inner sep=1pt] at (-101:24) {$2r$};
  \node[anchor=west, font=\small, text width=72mm, align=left] at (36,0) {%
  どちらの球面も、16本の線を全部受けとめる（本数は $r$ によらない）。\par\vspace{2pt}
  半径が2倍になると、球面の面積は4倍。同じ本数が4倍の面積に広がるので、\qty{1}{m^2} あたりの本数（電場の強さ）は $\tfrac14$。\par\vspace{2pt}
  太い弧は、同じ立体角の中の面。$2r$ では、たても横も2倍。};
\end{tikzpicture}
\bkfigcaption{点電荷から出る電気力線（実線）と、半径 $r$（実線の円）・$2r$（破線の円）の球面（切り口）。}\label{ph:電-2:fig-honsu}
\end{bkfigure}

\noindent 逆に考えると、電場が $1/r^2$ で弱まるのは、「決まった本数の線が、面積 $4\pi r^2$ の球面に広がって、薄まっていく」からだ、と言えます。これが、はてな？の答えの中心です（④でまとめます）。図\ref{ph:電-2:fig-honsu}の太い弧は、電荷から見て同じ広がり（\textbf{立体角}：8-8⑬）の中にある面です。距離が2倍になると、この面はたても横も2倍、面積は4倍になります。

\begin{phbaisuru}{距離 $r$}
球面の面積は4倍、線の本数は同じなので、\qty{1}{m^2} あたりの本数（電場）は $\tfrac14$ になります。電-1 のクーロンの法則の $1/r^2$ と同じです。\qty{1}{nC}なら、$r=\qty{0.1}{m}$ の球面（面積 \qty{0.126}{m^2}）に113本で、$113\div0.126\approx899$ 本/m$^2$。$r=\qty{0.2}{m}$（面積 \qty{0.503}{m^2}）では $113\div0.503\approx225$ 本/m$^2$。電-1 の表と同じ数です。
\end{phbaisuru}

\begin{motto}[4]
\phboxhead{窓を通る本数：立体角で数える}
電荷から見て、ある面がどれだけの広がりに見えるかを表すのが\textbf{立体角}でした（8-8⑬。単位 sr）。電荷の周り全体の立体角は $4\pi$ です。電気力線はどの向きにも同じように出ているので、立体角 $\Omega$ の中を通る本数は、全体の $\Omega/(4\pi)$ 倍です。
\[\text{通る本数}=\frac{Q}{\varepsilon_0}\times\frac{\Omega}{4\pi}\]
たとえば、電荷から見て半分の角が60°の円すい（$\Omega=\pi$、全体の $\tfrac14$：8-8⑬の表）の中には、\qty{1}{nC}なら $113\times\tfrac14\approx28$ 本の線が通ります。円すいの中の面なら、近くても遠くても、形が曲がっていても同じ28本です。②の「袋の形によらない」は、このことを袋全体（$\Omega=4\pi$）について言ったものです。
\end{motto}

%% ---------------------------------------------------------------
\bksubsection{2}{袋の形を変えても、本数は同じ}
\label{ph:電-2:sec2}
①では、点電荷を球面で囲みました。球でない袋で囲んだら、どうなるでしょう。図\ref{ph:電-2:fig-fukuro}のように、電荷を、でこぼこの袋で囲んでみます。

\begin{bkfigure}
\begin{tikzpicture}[x=1mm, y=1mm]
  \begin{scope}
  \foreach \a in {0,30,...,330} {\draw[sxc fl] (\a:2.4) -- (\a:25);}
  \draw[bkOrange!80!black, line width=1pt] plot[smooth cycle, tension=0.8] coordinates {(13,0) (10,9) (2,15) (-8,11) (-16,6) (-13,-6) (-6,-14) (6,-12) (9,-6)};
  \sxcplus{(0,0)}{2.4}
  \node[sxc lab, anchor=north] at (0,-26) {(ア) 電荷を囲む袋};
  \end{scope}
  \begin{scope}[shift={(66,0)}]
  \foreach \aa in {-21,-14,-7,0,7,14,21} {\draw[sxc fl] (-34,0) -- ++(\aa:58); \draw[sxc tip] ($(-34,0)+(\aa:53)$) -- ($(-34,0)+(\aa:54)$);}
  \draw[bkOrange!80!black, line width=1pt] plot[smooth cycle, tension=0.8] coordinates {(10,0) (7,9) (-1,13) (-9,9) (-12,1) (-8,-9) (2,-12) (8,-7)};
  \sxcplus{(-34,0)}{2.4}
  \node[sxc lab, anchor=south, fill=white, inner sep=1pt] at (-7,14) {入る線};
  \node[sxc lab, anchor=south, fill=white, inner sep=1pt] at (7,14) {出る線};
  \node[sxc lab, anchor=north] at (-4,-26) {(イ) 電荷が袋の外にあるとき};
  \end{scope}
\end{tikzpicture}
\bkfigcaption{袋を通りぬける電気力線。(ア) 袋の中の電荷から出た線は、袋の形がどうであっても、全部いつかは袋を通って外へ出る。(イ) 外の電荷の線は、袋に入った数だけ、反対側から出ていく。}\label{ph:電-2:fig-fukuro}
\end{bkfigure}

\noindent (ア) 電荷から出た線は、とちゅうで消えないので、袋の形がどうであっても、全部いつかは袋を通りぬけて外へ出ます。袋を通る本数は、やはり $Q/\varepsilon_0$ 本です。(イ) 電荷が袋の外にあるときは、線は袋に入っても、同じ線が必ず反対側から出ていきます。「出る本数」を正、「入る本数」を負として数えると、合わせて0本です。

\noindent このことから、次のきまりがわかります。
\begin{bkcenter}
\fcolorbox{bkTheme}{white}{\parbox{0.86\linewidth}{\centering 閉じた袋から出ていく電気力線の本数（出る本数－入る本数）は、\\袋の中の電荷の合計を $\varepsilon_0$ でわったものに等しい。袋の形にも、外の電荷にもよらない。}}
\end{bkcenter}
袋の中に電荷がいくつあっても、重ね合わせ（電-1 ③）で、それぞれの電荷の線をたせばよいので、「中の電荷の合計」で決まります。負の電荷は、線が入ってくるので、負の本数として数えます。8-8⑫では、同じことを、立方体の袋で流束の積分を計算して確かめました。

\Needspace{12\baselineskip}\phstage{te}{本数のきまりで、電場と電位を求める}{3}

\bksubsection{3}{ガウスの法則}
\label{ph:電-2:sec3}
②のきまりを、式で書きます。袋から出ていく電気力線の本数（\textbf{電場の流束}）を $\Phi_E$（ファイ・イー）、袋の中の電荷の合計を $Q_{\text{中}}$ とすると、
\pheq{\Phi_E=\frac{Q_{\text{中}}}{\varepsilon_0}}{N·m²/C ＝ C ÷ (C²/(N·m²))}
\label{ph:電-2:gauss}です。これを\textbf{ガウスの法則}といいます。\bkidx{がうすのほうそく}{ガウスの法則}\bkidx{でんばのりゅうそく}{電場の流束}ドイツの数学者・物理学者ガウスの名前がついています。

\phboxhead{本数を電場に直すこつ}
電場の強さは「\qty{1}{m^2}あたりの本数」でした。だから、袋の面の上で、
\begin{enumerate}[label={\bkcircled{\arabic*}}, leftmargin=2.4\zw, itemsep=1pt, topsep=2pt]
\item 電場が面に\textbf{垂直}で、強さ $E$ が面の上で\textbf{どこでも同じ}なら、通る本数は $E\times$（面積）。
\item 電場が面に\textbf{平行}なら、線は面をかすめるだけで、通る本数は0。
\end{enumerate}
袋を、この2つのどちらかの面だけでできるように選ぶと、$\Phi_E$ が「$E\times$面積」になり、ガウスの法則から $E$ がすぐに求まります。電荷の置き方が対称なとき、このような袋が見つかります（④）。

\begin{phflow}[ガウスの法則からクーロンの法則へ]
\phflowrow{$\Phi_E=\dfrac{Q_{\text{中}}}{\varepsilon_0}$}{点電荷を中心に\\半径 $r$ の球の袋}{$E\times4\pi r^2=\dfrac{Q}{\varepsilon_0}$}
\phflowrow{$E\times4\pi r^2=\dfrac{Q}{\varepsilon_0}$}{$4\pi r^2$ でわる\\$q$ をかける}{$F=k\,\dfrac{Qq}{r^2}$（クーロン）}
\end{phflow}

\begin{phyomi}[$\Phi_E=Q_{\text{中}}/\varepsilon_0$]
「袋から出ていく電気力線の本数は、中の電荷だけで決まる」。電荷は電気力線の「わき出し口」（負の電荷は「吸いこみ口」）で、線はそれ以外の所では生まれも消えもしない、ということです。8-9 の「わき出し」と同じ考え方です。
\end{phyomi}

\begin{seikaku}
ガウスの法則とクーロンの法則は、止まった電荷については同じ内容です（上の流れ図）。違いは、電荷が動くときに現れます。電荷が動くと、遠くの電場が変わるまでに時間がかかる（光の速さで伝わる）ので、クーロンの法則はそのままでは使えません。いっぽう、ガウスの法則は、電荷が動いていても成り立ちます。だから、電磁気学の基本の式（マクスウェル方程式：電-14）には、クーロンの法則ではなく、ガウスの法則が入っています。
\end{seikaku}

%% ---------------------------------------------------------------
\bksubsection{4}{対称な電荷：球・長い直線・広い平面}
\label{ph:電-2:sec4}
ガウスの法則を使って、対称な形に広がった電荷の電場を求めます。こつは、③の①②をみたす袋を、対称性から選ぶことです（8-11⑦の「袋を選ぶ3つの条件」と同じ。図\ref{ph:電-2:fig-taishou}）。


\phboxhead{(ア) 一様に帯電した球}
半径 $R$ の球の中に、電荷 $Q$ が一様に（どこでも同じ密度で）広がっています。対称性から、電場は中心から放射状で、中心から同じ距離の所では同じ強さです。半径 $r$ の球の袋を選びます。
\begin{itemize}[leftmargin=1.2\zw, itemsep=2pt, topsep=2pt, label={\color{bkTheme}●}]
\item \textbf{球の外}（$r\geqq R$）：袋の中の電荷は $Q$ 全部なので、$E\times4\pi r^2=Q/\varepsilon_0$ から $E=kQ/r^2$。\textbf{球の外では、全部の電荷が中心の1点に集まっているときと同じ}です。
\item \textbf{球の中}（$r<R$）：袋の中の電荷は、体積の割合で $Q\times r^3/R^3$ です。
\[E\times4\pi r^2=\frac{Q}{\varepsilon_0}\times\frac{r^3}{R^3}\qquad\text{から}\qquad E=\frac{kQ}{R^3}\,r\]
で、中心からの距離に\textbf{比例}します。中心では0です。
\end{itemize}
これは、8-11⑦で、地球の重力について同じ計算をしたのと同じです。

\begin{bkfigure}
\begin{tikzpicture}[x=1mm, y=1mm]
  \begin{scope}[shift={(18,0)}]
    \fill[bkRed!15] (0,0) circle (8);
    \draw[bkRed!70!black, line width=0.6pt] (0,0) circle (8);
    \foreach \a in {20,80,...,320} {\sxcp{(\a:4.5)}}
    \draw[bkOrange!80!black, line width=1pt, dash pattern=on 2.4pt off 1.4pt] (0,0) circle (14);
    \foreach \a in {0,45,...,315} {\draw[sxc arr] (\a:14.6) -- (\a:19.5);}
    \node[sxc lab, anchor=north, align=center] at (0,-21) {(ア) 球\\→ 同じ中心の球の袋};
  \end{scope}
  \begin{scope}[shift={(64,0)}]
    \draw[bkRed!70!black, line width=2.2pt] (0,-18) -- (0,18);
    \foreach \yy in {-15,-9,...,15} {\sxcp{(-2.2,\yy)}}
    \draw[bkOrange!80!black, line width=1pt, dash pattern=on 2.4pt off 1.4pt] (-9,-10) rectangle (9,10);
    \draw[bkOrange!80!black, line width=0.6pt, dash pattern=on 2.4pt off 1.4pt] (0,10) ellipse (9 and 2.4) (0,-10) ellipse (9 and 2.4);
    \foreach \yy in {-5,0,5} {\draw[sxc arr] (9.6,\yy) -- (15,\yy); \draw[sxc arr] (-9.6,\yy) -- (-15,\yy);}
    \node[sxc lab, anchor=north, align=center] at (0,-21) {(イ) 長い直線\\→ 同じ軸の円柱の袋};
  \end{scope}
  \begin{scope}[shift={(108,0)}]
    \fill[bkRed!15] (-16,-1) rectangle (16,1);
    \draw[bkRed!70!black, line width=0.6pt] (-16,-1) rectangle (16,1);
    \draw[bkOrange!80!black, line width=1pt, dash pattern=on 2.4pt off 1.4pt] (-6,-8) rectangle (6,8);
    \draw[sxc arr] (0,8.6) -- (0,15);
    \draw[sxc arr] (0,-8.6) -- (0,-15);
    \node[sxc lab, anchor=north, align=center] at (0,-21) {(ウ) 広い平面\\→ はさむ箱の袋};
  \end{scope}
\end{tikzpicture}
\bkfigcaption{対称な電荷と、選ぶ袋（破線）。矢印は電場。(ア) 球面の上で電場は垂直で同じ強さ。(イ) 円柱の側面で垂直、上下の面では平行。(ウ) 箱の上下の面で垂直、横の面では平行。}\label{ph:電-2:fig-taishou}
\end{bkfigure}

\begin{bkfigure}
\begin{tikzpicture}
\begin{axis}[bkplot, width=110mm, height=46mm, xmin=0, xmax=0.33, ymin=0, ymax=10500, xtick={0.05,0.1,0.2,0.3}, xticklabel style={/pgf/number format/fixed}, ytick={2248,4495,8990}, yticklabels={2248,4495,8990}, xlabel={$r$（m）}, ylabel={$E$（N/C）}]
  \addplot[bkcurve, domain=0:0.1, samples=10] {89900*x};
  \addplot[bkcurve, domain=0.1:0.32, samples=80] {89.9/x^2};
  \draw[bkGray, line width=0.5pt, dashed] (axis cs:0.1,0) -- (axis cs:0.1,10000);
  \fill[bkRed] (axis cs:0.05,4495) circle (1.8pt) (axis cs:0.2,2247.5) circle (1.8pt);
  \node[font=\sffamily\footnotesize, anchor=south east, fill=white, inner sep=1pt] at (axis cs:0.095,8000) {球の中：$r$ に比例};
  \node[font=\sffamily\footnotesize, anchor=south west, fill=white, inner sep=1pt] at (axis cs:0.13,5600) {球の外：$1/r^2$};
\end{axis}
\end{tikzpicture}
\bkfigcaption{一様に帯電した球（$Q=\qty{1e-8}{C}$、半径 $R=\qty{0.1}{m}$）の電場。表面（破線）でいちばん強い。●は例題1の2点。}\label{ph:電-2:fig-kyu}
\end{bkfigure}

\phboxhead{(イ) 長い直線}
1 m あたり $\lambda$ の電荷が一様にある長い直線を、半径 $r$・長さ $l$ の円柱の袋で囲みます。電場は直線から放射状で、側面では垂直・同じ強さ、上下の面では平行（通る本数0）です。袋の中の電荷は $\lambda l$ なので、
\[E\times2\pi rl=\frac{\lambda l}{\varepsilon_0}\qquad\text{より}\qquad E=\frac{\lambda}{2\pi\varepsilon_0r}=\frac{2k\lambda}{r}\]
です。電-1 ⑦で、積分を使って求めた答えと同じです（例題2）。

\phboxhead{(ウ) 広い平面}
1 m$^2$ あたり $\sigma$ の電荷が一様にある広い平面を、上下の面の面積が $S$ の箱ではさみます。電場は面に垂直で、箱の上下の面を外向きに通り、横の面では平行です。袋の中の電荷は $\sigma S$ なので、
\[E\times2S=\frac{\sigma S}{\varepsilon_0}\qquad\text{より}\qquad E=\frac{\sigma}{2\varepsilon_0}\]
で、電-1 ⑦と同じく、面からの距離によりません。

\begin{phflow}[ガウスの法則と対称性]
\phflowrow{$\Phi_E=\dfrac{Q_{\text{中}}}{\varepsilon_0}$}{一様な球の外\\球の袋}{$E=k\,\dfrac{Q}{r^2}$}
\phflowrow{$\Phi_E=\dfrac{Q_{\text{中}}}{\varepsilon_0}$}{長い直線\\円柱の袋}{$E=\dfrac{\lambda}{2\pi\varepsilon_0r}$}
\phflowrow{$\Phi_E=\dfrac{Q_{\text{中}}}{\varepsilon_0}$}{広い平面\\はさむ箱の袋}{$E=\dfrac{\sigma}{2\varepsilon_0}$}
\end{phflow}

\phboxhead{はてな？の答え}
\label{ph:電-2:hatena}\bkhatenaref の答えです。電荷から出る電気力線の本数は決まっていて（$Q/\varepsilon_0$ 本）、とちゅうで増えも減りもしません。その線が、距離 $r$ の所では面積 $4\pi r^2$ の球面全体に広がっています。電場の強さは「\qty{1}{m^2}あたりの本数」なので、
\[E=\frac{Q/\varepsilon_0}{4\pi r^2}\propto\frac{1}{r^2}\]
です。電気の力が\textbf{距離の2乗}で弱まるのは、\textbf{わたしたちの空間が3次元で、球面の面積が $r^2$ に比例するから}なのです。音や光の強さが距離の2乗に反比例して弱まる（波-7 ②の逆2乗の法則）のも、まったく同じ理由です。やってみようのライトの光の明るさが、距離2倍で $\tfrac14$ になったことを思い出しましょう。

この考えで、電荷の形による違いも説明できます。長い直線の電荷では、線は円柱の側面（面積が $r$ に比例）に広がるので $1/r$、広い平面では、線は広がらずにまっすぐ進むので、距離によりません。もし空間が2次元（平面だけの世界）なら、点電荷からの線は円周（長さ $2\pi r$）に広がるので、電気の力は $1/r$ で弱まるはずです。

%% ---------------------------------------------------------------
\Needspace{12\baselineskip}\bksubsection{5}{電位：電気の「高さ」}
\label{ph:電-2:sec5}
重力の中で物を持ち上げると、その物は位置エネルギー $mgh$ をたくわえます（力-4）。電場の中でも、電荷を、電場の力にさからって運ぶと、仕事をした分だけ\textbf{電気の位置エネルギー}がたまります。

\phboxhead{点電荷の周りの位置エネルギー}
正の点電荷 $Q$ から十分遠い所（「無限に遠い所」と考えます）から、正の電荷 $q$ を、$Q$ から距離 $r$ の所まで、ゆっくり運んでくるとします。$q$ ははね返す力 $F=kQq/x^2$（$x$ は $Q$ からの距離）を受けるので、それにさからう力で押していく仕事が要ります。力が距離で変わるので、仕事は「力のグラフの面積」です（図\ref{ph:電-2:fig-menseki}）。
\[W=\int_r^{\infty}k\,\frac{Qq}{x^2}\,\dd x=kQq\Bigl[-\frac1x\Bigr]_r^{\infty}=k\,\frac{Qq}{r}\]
この仕事が、$q$ の位置エネルギー $U$ になります。無限に遠い所を位置エネルギー0の基準にすると、
\pheq{U=k\,\frac{Qq}{r}}{J ＝ $\dfrac{\text{N·m}^2}{\text{C}^2}$×$\dfrac{\text{C}^2}{\text{m}}$}

\begin{bkfigure}
\begin{tikzpicture}
\begin{axis}[bkplot, width=112mm, height=46mm, xmin=0, xmax=0.52, ymin=0, ymax=2000, xtick={0.1,0.2,0.3,0.4,0.5}, xticklabel style={/pgf/number format/fixed}, ytick={899}, xlabel={$x$（m）}, ylabel={$E$（N/C）}]
  \addplot[draw=none, fill=bkOrange!25, domain=0.1:0.5, samples=60] {8.99/x^2} \closedcycle;
  \addplot[bkcurve, domain=0.068:0.5, samples=80] {8.99/x^2};
  \draw[bkGray, line width=0.5pt, dashed] (axis cs:0.1,0) -- (axis cs:0.1,899);
  \node[font=\sffamily\footnotesize, anchor=west, fill=white, inner sep=1pt] at (axis cs:0.18,900) {色の部分の面積＝$\phi$（$r=0.1$ から先）};
  \node[font=\sffamily\footnotesize, anchor=south, fill=white, inner sep=1pt] at (axis cs:0.1,1300) {$r$};
\end{axis}
\end{tikzpicture}
\bkfigcaption{\qty{1}{nC}の点電荷の電場のグラフ。$r=\qty{0.1}{m}$ から遠くまでの面積（色の部分。右へずっと続く）が、$r$ の所の電位 $\phi=kQ/r\approx\qty{89.9}{V}$ になる。}\label{ph:電-2:fig-menseki}
\end{bkfigure}

\phboxhead{電位：1 C あたりの位置エネルギー}
位置エネルギー $U$ は、運ぶ電荷 $q$ に比例します。そこで、$q$ でわって「\qty{1}{C}あたりの位置エネルギー」にしたものを、その場所の\textbf{電位}といい、$\phi$ と書きます。\bkidx{でんい}{電位}
\pheq{\phi=\frac{U}{q},\qquad\text{点電荷の周りでは}\quad\phi=k\,\frac{Q}{r}}{V ＝ J/C}
\label{ph:電-2:deni}電位の単位は J/C で、これを\textbf{ボルト}〔V〕といいます。電場が「1\,C あたりの力」だったのと同じく、電位は「1\,C あたりの位置エネルギー」で、置く電荷によらず、場所だけで決まります。電荷 $q$ を電位 $\phi$ の所に置くと、位置エネルギーは
\[U=q\phi\]
です。電場の単位 N/C は、V/m とも書けます（N/C＝N·m/(C·m)＝J/(C·m)＝V/m）。

\begin{phyomi}[$\phi=k\,\dfrac{Q}{r}$]
「正の電荷の近くは電位が高い山、負の電荷の近くは深い谷」。電位は、電場と違って\textbf{向きのない数}です。いくつかの電荷の電位は、それぞれの電位を\textbf{ふつうにたす}だけで求められます（矢印のたし算はいらない）。距離の1乗に反比例するので、電場（2乗）よりゆっくり小さくなります。
\end{phyomi}

\begin{phbaisuru}{距離 $r$}
電位は $\tfrac12$ になります（電場は $\tfrac14$）。\qty{1}{nC}なら、\qty{0.1}{m}で\qty{89.9}{V}、\qty{0.2}{m}で\qty{45.0}{V}です。電位が $1/r$ で、電場が $1/r^2$ なのは、電位が「電場のグラフの面積」だからです（$1/x^2$ を積分すると $1/x$）。
\par\smallskip\noindent\begin{tikzpicture}
\begin{axis}[bkplot, width=100mm, height=44mm, xmin=0, xmax=0.55, ymin=0, ymax=1.15, xtick={0.1,0.2,0.3,0.4,0.5}, xticklabel style={/pgf/number format/fixed}, ytick={0.25,0.5,1}, xlabel={$r$（m）}, ylabel={\qty{0.1}{m}の値を1とした大きさ}]
  \addplot[bkcurve, domain=0.087:0.52, samples=80] {0.1/x};
  \addplot[bkOrange, line width=1.6pt, dash pattern=on 4pt off 2pt, domain=0.095:0.52, samples=80] {0.01/x^2};
  \fill[bkRed] (axis cs:0.2,0.5) circle (1.8pt) (axis cs:0.2,0.25) circle (1.8pt);
  \node[font=\sffamily\footnotesize, anchor=south west, fill=white, inner sep=1pt] at (axis cs:0.22,0.5) {電位 $\phi$（実線）：半分};
  \node[font=\sffamily\footnotesize, anchor=north west, fill=white, inner sep=1pt] at (axis cs:0.25,0.2) {電場 $E$（破線）：$\tfrac14$};
\end{axis}
\end{tikzpicture}
\end{phbaisuru}

\phboxhead{電圧と仕事}
2点 A・B の電位の差を、A と B のあいだの\textbf{電圧}（電位差）といい、$V$ と書きます。\bkidx{でんあつ}{電圧}電荷 $q$ が A から B へ動くとき、電場の力がする仕事は、位置エネルギーの減った分なので、
\pheq{W=q(\phi_{\mathrm A}-\phi_{\mathrm B})=qV}{J ＝ C × V}
です。電場の力は保存力なので、この仕事は、A から B へどんな道を通っても同じです（力-4、8-7⑨）。乾電池の「\qty{1.5}{V}」は、電池の＋極と－極のあいだの電圧で、\qty{1}{C}の電荷が＋極から回路を通って－極へもどるあいだに、\qty{1.5}{J}の仕事ができる、という意味です（電-4）。

\begin{phmijika}[電位は水の高さ、電圧は高さの差]
電位を「水の高さ」、正の電荷を「水」にたとえると、わかりやすくなります。水は高い所から低い所へ流れ、そのとき、高さの差に比例した仕事（水車を回すなど）をします。正の電荷も、電位の高い所から低い所へ動こうとし、電位の差（電圧）に比例した仕事をします。乾電池は、水を高い所へくみ上げるポンプのようなもので、電荷を\qty{1.5}{V}だけ高い電位へ持ち上げます（電-4）。
\par\smallskip\noindent\begin{tikzpicture}[x=1mm, y=1mm]
  \fill[bkBlue!15] (0,20) rectangle (16,26);
  \fill[bkBlue!15] (40,0) rectangle (56,6);
  \draw[bkInk, line width=0.6pt] (0,28) -- (0,20) -- (16,20) -- (16,28) (40,8) -- (40,0) -- (56,0) -- (56,8);
  \draw[bkBlue!70!black, line width=1.6pt, -{Stealth[length=2.6mm]}] (16,22) .. controls (28,22) and (28,4) .. (40,4);
  \draw[sxc dim] (60,3) -- (60,23);
  \node[sxc lab, anchor=west] at (61,13) {高さの差};
  \node[sxc lab, anchor=south] at (8,27) {高い所};
  \node[sxc lab, anchor=north] at (48,-1) {低い所};
  \begin{scope}[shift={(84,0)}]
  \fill[bkRed!15] (0,24) rectangle (16,26);
  \fill[bkBlue!15] (0,0) rectangle (16,2);
  \draw[bkRed!70!black, line width=0.6pt] (0,24) rectangle (16,26);
  \draw[bkBlue!70!black, line width=0.6pt] (0,0) rectangle (16,2);
  \sxcplus{(8,18)}{2}
  \draw[bkRed!80!black, line width=1.2pt, -{Stealth[length=2.4mm]}] (8,15.6) -- (8,6);
  \draw[sxc dim] (20,1) -- (20,25);
  \node[sxc lab, anchor=west] at (21,13) {電圧 $V$};
  \node[sxc lab, anchor=south] at (8,26.5) {高い電位};
  \node[sxc lab, anchor=north] at (8,-1) {低い電位};
  \end{scope}
\end{tikzpicture}
\par\noindent（左：水は高い所から低い所へ流れる。右：正の電荷は電位の高い所から低い所へ動き、電場は $qV$ の仕事をする。）
\end{phmijika}

\begin{seikaku}
電位の0の基準は、自由に決めてよい数です。高さの基準を海面にしても、床にしてもよいのと同じです。物理で大切なのは電位の\textbf{差}だけなので、点電荷の周りでは「無限に遠い所を0」、回路では「地面（アース）を0」とするのがふつうです。数学 8-12⑦ の点電荷の電位 $\phi=\dfrac{Q}{4\pi\varepsilon_0}\times\dfrac1r$ は、この $\phi=kQ/r$ と同じものです。
\end{seikaku}

%% ---------------------------------------------------------------
\bksubsection{6}{一様な電場：$V=Ed$ と電子ボルト}
\label{ph:電-2:sec6}
向かい合った広い2枚の板のあいだの、一様な電場 $E$ を考えます（電-1 ⑦）。正の電荷 $q$ を、電場の向きに距離 $d$ だけ運ぶと、電場の力 $qE$ がする仕事は $qEd$ です。これが $qV$ に等しいので、
\pheq{V=Ed}{V ＝ V/m × m}
\label{ph:電-2:ved}です。間かく\qty{1}{cm}の板に\qty{100}{V}の電圧をかけると、あいだの電場は $100\div0.01$ で、\qty{1.0e4}{V/m}です（8-4⑧ と同じ数）。電場の単位を V/m と書くと、「1 m あたり何 V 下がるか」と読めて便利です。

\begin{bkfigure}
\begin{tikzpicture}[x=1mm, y=1mm]
  \fill[bkRed!15] (0,30) rectangle (60,32);
  \fill[bkBlue!15] (0,-2) rectangle (60,0);
  \draw[bkRed!70!black, line width=0.6pt] (0,30) rectangle (60,32);
  \draw[bkBlue!70!black, line width=0.6pt] (0,-2) rectangle (60,0);
  \foreach \xx in {4,12,...,60} {\sxcp{(\xx,33.8)}\sxcm{(\xx,-3.8)}}
  \foreach \yy/\vv in {7.5/25,15/50,22.5/75} {\draw[sxc eq] (0,\yy) -- (60,\yy); \node[sxc lab, anchor=west, text=bkGreen!40!black] at (61,\yy) {\qty{\vv}{V}};}
  \foreach \xx in {8,22,38,52} {\draw[sxc arr, line width=1.1pt] (\xx,28.5) -- (\xx,1.5);}
  \node[sxc lab, anchor=west] at (61,31) {\qty{100}{V}};
  \node[sxc lab, anchor=west] at (61,-1) {\qty{0}{V}};
  \draw[sxc dim] (-4,0) -- (-4,30);
  \node[sxc lab, anchor=east] at (-5,15) {$d$};
  \node[anchor=west, font=\small, text width=44mm, align=left] at (78,15) {%
  実線の矢印：電場（下向き、どこでも同じ）。\par\vspace{2pt}
  破線：等電位面。$\tfrac{d}{4}$ 下がるごとに\qty{25}{V}ずつ下がる。};
\end{tikzpicture}
\bkfigcaption{一様な電場の中の電位。電場の向きに進むと、電位は一定の割合で下がる（$V=Ed$）。}\label{ph:電-2:fig-ichiyou}
\end{bkfigure}

\phboxhead{電子ボルト}
電荷 $e$ の粒（電子・陽子）が、電圧 $V$ で加速されると、運動エネルギーは $eV$ だけ増えます（エネルギーの保存）。たとえば\qty{1}{V}で加速された電子は、$1.60\times10^{-19}$\,J のエネルギーを得ます。このエネルギーを\textbf{1電子ボルト}〔eV〕といいます。\bkidx{でんしぼると}{電子ボルト}
\pheq{\qty{1}{eV}=\qty{1.60e-19}{J}}{J ＝ C × V}
原子や分子、電子の世界では、J ではなく eV でエネルギーを表すと、ちょうどよい大きさの数になります。電圧 $V$ で、止まっていた電子（質量 $m_{\mathrm e}$）を加速すると、$\tfrac12m_{\mathrm e}v^2=eV$ から
\[v=\sqrt{\frac{2eV}{m_{\mathrm e}}}\]
です。\qty{100}{V}で約 \qty{5.9e6}{m/s}、\qty{1}{kV}で約 \qty{1.9e7}{m/s}（光の速さの約6\,\%）になります（例題5）。

\begin{phmijika}[電子を電圧で加速する]
昔のテレビ（ブラウン管）は、電子を数万 V で加速して、画面の内側の蛍光体に当てて光らせていました。電子顕微鏡や、病院の X 線の装置も、電子を高い電圧で加速しています。X 線の写真をとる装置では、数万〜十数万 V で加速した電子を金属に当てて X 線を出すので、X 線のエネルギーも数十 keV（キロ電子ボルト）になります。
\end{phmijika}

\begin{phbaisuru}{電圧 $V$}
電子が得るエネルギー $eV$ は2倍、速さは $\sqrt2\approx1.41$ 倍です（速さは電圧の平方根に比例）。速さを2倍にするには、電圧を4倍にしなければなりません（運動エネルギーは速さの2乗）。
\end{phbaisuru}

%% ---------------------------------------------------------------
\Needspace{14\baselineskip}\bksubsection{7}{等電位面と電気力線}
\label{ph:電-2:sec7}
電位が同じ点をつないだ面を\textbf{等電位面}（平面の図では等電位線）といいます。\bkidx{とうでんいめん}{等電位面}地図の等高線と同じで、電位が「高さ」、等電位面が「等高線」です。

\begin{bkfigure}
\begin{tikzpicture}[x=0.6mm, y=0.6mm]
\begin{scope}
  \clip (-52,-30) rectangle (52,30);
\draw[sxc eq] plot[smooth, tension=0.6] coordinates {(-9.7,0.0) (-9.9,1.2) (-10.3,2.4) (-11.0,3.5) (-11.9,4.3) (-12.9,5.0) (-14.1,5.4) (-15.4,5.5) (-16.6,5.3) (-17.8,4.8) (-18.8,4.1) (-19.6,3.2) (-20.2,2.1) (-20.5,0.9) (-20.6,-0.4) (-20.4,-1.6) (-19.9,-2.7) (-19.1,-3.7) (-18.2,-4.6) (-17.1,-5.1) (-15.9,-5.4) (-14.6,-5.4) (-13.4,-5.2) (-12.3,-4.6) (-11.3,-3.9) (-10.6,-2.9) (-10.0,-1.7) (-9.8,-0.5) (-9.7,0.0)};
\draw[sxc eq] plot[smooth, tension=0.6] coordinates {(-6.4,0.0) (-6.5,1.2) (-6.7,2.5) (-7.1,3.7) (-7.6,4.8) (-8.2,5.9) (-9.0,6.9) (-9.9,7.7) (-10.9,8.5) (-12.0,9.0) (-13.2,9.5) (-14.4,9.8) (-15.7,9.9) (-16.9,9.8) (-18.1,9.6) (-19.3,9.2) (-20.5,8.7) (-21.5,8.0) (-22.5,7.2) (-23.3,6.3) (-24.1,5.3) (-24.6,4.2) (-25.1,3.0) (-25.4,1.8) (-25.5,0.6) (-25.5,-0.7) (-25.4,-1.9) (-25.1,-3.1) (-24.6,-4.3) (-24.0,-5.4) (-23.3,-6.4) (-22.4,-7.3) (-21.4,-8.1) (-20.4,-8.7) (-19.2,-9.3) (-18.0,-9.6) (-16.8,-9.8) (-15.6,-9.9) (-14.3,-9.7) (-13.1,-9.5) (-11.9,-9.0) (-10.8,-8.4) (-9.8,-7.6) (-9.0,-6.8) (-8.2,-5.8) (-7.5,-4.7) (-7.0,-3.6) (-6.7,-2.4) (-6.5,-1.1) (-6.4,0.0)};
\draw[sxc eq] plot[smooth, tension=0.6] coordinates {(-3.2,0.0) (-3.3,1.2) (-3.4,2.5) (-3.5,3.7) (-3.8,5.0) (-4.1,6.2) (-4.5,7.4) (-5.0,8.5) (-5.5,9.6) (-6.2,10.7) (-6.9,11.7) (-7.7,12.7) (-8.6,13.5) (-9.6,14.3) (-10.6,15.1) (-11.7,15.7) (-12.8,16.2) (-14.0,16.7) (-15.2,17.0) (-16.4,17.2) (-17.6,17.3) (-18.9,17.3) (-20.1,17.3) (-21.4,17.1) (-22.6,16.8) (-23.8,16.4) (-24.9,15.9) (-26.1,15.4) (-27.1,14.7) (-28.2,14.0) (-29.1,13.2) (-30.0,12.4) (-30.9,11.4) (-31.6,10.4) (-32.3,9.4) (-32.9,8.3) (-33.5,7.2) (-34.0,6.0) (-34.3,4.8) (-34.6,3.6) (-34.8,2.4) (-35.0,1.2) (-35.0,-0.1) (-34.9,-1.3) (-34.8,-2.6) (-34.6,-3.8) (-34.3,-5.0) (-33.9,-6.2) (-33.4,-7.4) (-32.9,-8.5) (-32.2,-9.6) (-31.5,-10.6) (-30.7,-11.6) (-29.9,-12.5) (-29.0,-13.3) (-28.0,-14.1) (-27.0,-14.8) (-25.9,-15.5) (-24.8,-16.0) (-23.6,-16.5) (-22.4,-16.8) (-21.2,-17.1) (-19.9,-17.3) (-18.7,-17.3) (-17.5,-17.3) (-16.2,-17.2) (-15.0,-16.9) (-13.8,-16.6) (-12.6,-16.1) (-11.5,-15.6) (-10.4,-15.0) (-9.4,-14.2) (-8.5,-13.4) (-7.6,-12.5) (-6.8,-11.6) (-6.1,-10.5) (-5.4,-9.5) (-4.9,-8.3) (-4.4,-7.2) (-4.0,-6.0) (-3.7,-4.8) (-3.5,-3.5) (-3.3,-2.3) (-3.2,-1.1) (-3.2,0.0)};
\draw[sxc eq] plot[smooth, tension=0.6] coordinates {(9.7,0.0) (9.9,1.2) (10.3,2.4) (11.0,3.5) (11.9,4.3) (12.9,5.0) (14.1,5.4) (15.4,5.5) (16.6,5.3) (17.8,4.8) (18.8,4.1) (19.6,3.2) (20.2,2.1) (20.5,0.9) (20.6,-0.4) (20.4,-1.6) (19.9,-2.7) (19.1,-3.7) (18.2,-4.6) (17.1,-5.1) (15.9,-5.4) (14.6,-5.4) (13.4,-5.2) (12.3,-4.6) (11.3,-3.9) (10.6,-2.9) (10.0,-1.7) (9.8,-0.5) (9.7,0.0)};
\draw[sxc eq] plot[smooth, tension=0.6] coordinates {(6.4,0.0) (6.5,1.2) (6.7,2.5) (7.1,3.7) (7.6,4.8) (8.2,5.9) (9.0,6.9) (9.9,7.7) (10.9,8.5) (12.0,9.0) (13.2,9.5) (14.4,9.8) (15.7,9.9) (16.9,9.8) (18.1,9.6) (19.3,9.2) (20.5,8.7) (21.5,8.0) (22.5,7.2) (23.3,6.3) (24.1,5.3) (24.6,4.2) (25.1,3.0) (25.4,1.8) (25.5,0.6) (25.5,-0.7) (25.4,-1.9) (25.1,-3.1) (24.6,-4.3) (24.0,-5.4) (23.3,-6.4) (22.4,-7.3) (21.4,-8.1) (20.4,-8.7) (19.2,-9.3) (18.0,-9.6) (16.8,-9.8) (15.6,-9.9) (14.3,-9.7) (13.1,-9.5) (11.9,-9.0) (10.8,-8.4) (9.8,-7.6) (9.0,-6.8) (8.2,-5.8) (7.5,-4.7) (7.0,-3.6) (6.7,-2.4) (6.5,-1.1) (6.4,0.0)};
\draw[sxc eq] plot[smooth, tension=0.6] coordinates {(3.2,0.0) (3.3,1.2) (3.4,2.5) (3.5,3.7) (3.8,5.0) (4.1,6.2) (4.5,7.4) (5.0,8.5) (5.5,9.6) (6.2,10.7) (6.9,11.7) (7.7,12.7) (8.6,13.5) (9.6,14.3) (10.6,15.1) (11.7,15.7) (12.8,16.2) (14.0,16.7) (15.2,17.0) (16.4,17.2) (17.6,17.3) (18.9,17.3) (20.1,17.3) (21.4,17.1) (22.6,16.8) (23.8,16.4) (24.9,15.9) (26.1,15.4) (27.1,14.7) (28.2,14.0) (29.1,13.2) (30.0,12.4) (30.9,11.4) (31.6,10.4) (32.3,9.4) (32.9,8.3) (33.5,7.2) (34.0,6.0) (34.3,4.8) (34.6,3.6) (34.8,2.4) (35.0,1.2) (35.0,-0.1) (34.9,-1.3) (34.8,-2.6) (34.6,-3.8) (34.3,-5.0) (33.9,-6.2) (33.4,-7.4) (32.9,-8.5) (32.2,-9.6) (31.5,-10.6) (30.7,-11.6) (29.9,-12.5) (29.0,-13.3) (28.0,-14.1) (27.0,-14.8) (25.9,-15.5) (24.8,-16.0) (23.6,-16.5) (22.4,-16.8) (21.2,-17.1) (19.9,-17.3) (18.7,-17.3) (17.5,-17.3) (16.2,-17.2) (15.0,-16.9) (13.8,-16.6) (12.6,-16.1) (11.5,-15.6) (10.4,-15.0) (9.4,-14.2) (8.5,-13.4) (7.6,-12.5) (6.8,-11.6) (6.1,-10.5) (5.4,-9.5) (4.9,-8.3) (4.4,-7.2) (4.0,-6.0) (3.7,-4.8) (3.5,-3.5) (3.3,-2.3) (3.2,-1.1) (3.2,0.0)};
\draw[sxc eq] (0,-30) -- (0,30);
\draw[sxc fl] plot[smooth] coordinates {(-15.0,0.0) (-12.8,0.0) (-11.6,0.0) (-10.3,0.0) (-9.1,0.0) (-7.8,0.0) (-6.6,0.0) (-5.3,0.0) (-4.1,0.0) (-2.8,0.0) (-1.6,0.0) (-0.3,0.0) (0.9,0.0) (2.2,0.0) (3.4,0.0) (4.7,0.0) (5.9,0.0) (7.2,0.0) (8.4,0.0) (9.7,0.0) (10.9,0.0) (12.2,0.0) (15.0,0.0)}; \draw[sxc tip] (-0.60,0.00) -- (-0.10,0.00);
\draw[sxc fl] plot[smooth] coordinates {(-15.0,0.0) (-13.3,1.4) (-12.4,2.2) (-11.4,3.0) (-10.4,3.8) (-9.4,4.5) (-8.3,5.2) (-7.3,5.8) (-6.2,6.4) (-5.0,7.0) (-3.9,7.4) (-2.7,7.7) (-1.4,7.9) (-0.2,8.0) (1.1,8.0) (2.3,7.8) (3.5,7.5) (4.7,7.1) (5.9,6.6) (7.0,6.0) (8.0,5.4) (9.1,4.7) (10.1,4.0) (11.1,3.2) (12.1,2.4) (13.0,1.6) (15.0,0.0)}; \draw[sxc tip] (-0.50,8.00) -- (0.00,8.00);
\draw[sxc fl] plot[smooth] coordinates {(-15.0,0.0) (-14.6,2.2) (-14.4,3.4) (-14.1,4.6) (-13.9,5.8) (-13.6,7.1) (-13.2,8.3) (-12.8,9.4) (-12.4,10.6) (-11.9,11.8) (-11.3,12.9) (-10.7,14.0) (-10.0,15.0) (-9.3,16.0) (-8.5,16.9) (-7.5,17.8) (-6.6,18.6) (-5.5,19.3) (-4.4,19.8) (-3.2,20.3) (-2.0,20.6) (-0.8,20.7) (0.5,20.8) (1.7,20.6) (2.9,20.4) (4.1,20.0) (5.2,19.4) (6.3,18.8) (7.3,18.0) (8.2,17.2) (9.1,16.3) (9.9,15.3) (10.6,14.3) (11.2,13.2) (11.8,12.1) (12.3,10.9) (12.7,9.8) (13.1,8.6) (13.5,7.4) (13.8,6.2) (14.1,4.9) (14.3,3.7) (14.6,2.5) (15.0,0.0)}; \draw[sxc tip] (0.20,20.85) -- (0.70,20.77);
\draw[sxc fl] plot[smooth] coordinates {(-15.0,0.0) (-16.1,1.9) (-16.7,3.0) (-17.3,4.1) (-17.9,5.2) (-18.5,6.3) (-19.0,7.4) (-19.6,8.5) (-20.1,9.7) (-20.6,10.8) (-21.1,12.0) (-21.5,13.2) (-21.9,14.3) (-22.3,15.5) (-22.6,16.7) (-23.0,17.9) (-23.3,19.2) (-23.5,20.4) (-23.7,21.6) (-23.9,22.8) (-24.1,24.1) (-24.2,25.3) (-24.3,26.6) (-24.3,27.8) (-24.4,29.1) (-24.3,29.8)}; \draw[sxc tip] (-21.40,12.92) -- (-21.57,13.39);
\draw[sxc fl] plot[smooth] coordinates {(-15.0,0.0) (-17.1,0.8) (-18.2,1.2) (-19.4,1.6) (-20.6,2.1) (-21.8,2.5) (-22.9,3.0) (-24.1,3.4) (-25.2,3.9) (-26.4,4.3) (-27.6,4.8) (-28.7,5.3) (-29.9,5.8) (-31.0,6.3) (-32.1,6.8) (-33.3,7.3) (-34.4,7.8) (-35.6,8.3) (-36.7,8.9) (-37.8,9.4) (-38.9,10.0) (-40.1,10.5) (-41.2,11.1) (-42.3,11.6) (-43.4,12.2) (-44.5,12.8) (-45.6,13.4) (-46.7,14.0) (-47.8,14.6) (-48.9,15.2) (-50.0,15.8) (-51.1,16.4) (-52.0,16.9)}; \draw[sxc tip] (-33.03,7.18) -- (-33.48,7.38);
\draw[sxc fl] plot[smooth] coordinates {(-15.0,0.0) (-17.1,-0.8) (-18.2,-1.2) (-19.4,-1.6) (-20.6,-2.1) (-21.8,-2.5) (-22.9,-3.0) (-24.1,-3.4) (-25.2,-3.9) (-26.4,-4.3) (-27.6,-4.8) (-28.7,-5.3) (-29.9,-5.8) (-31.0,-6.3) (-32.1,-6.8) (-33.3,-7.3) (-34.4,-7.8) (-35.6,-8.3) (-36.7,-8.9) (-37.8,-9.4) (-38.9,-10.0) (-40.1,-10.5) (-41.2,-11.1) (-42.3,-11.6) (-43.4,-12.2) (-44.5,-12.8) (-45.6,-13.4) (-46.7,-14.0) (-47.8,-14.6) (-48.9,-15.2) (-50.0,-15.8) (-51.1,-16.4) (-52.0,-16.9)}; \draw[sxc tip] (-33.03,-7.18) -- (-33.48,-7.38);
\draw[sxc fl] plot[smooth] coordinates {(-15.0,0.0) (-16.1,-1.9) (-16.7,-3.0) (-17.3,-4.1) (-17.9,-5.2) (-18.5,-6.3) (-19.0,-7.4) (-19.6,-8.5) (-20.1,-9.7) (-20.6,-10.8) (-21.1,-12.0) (-21.5,-13.2) (-21.9,-14.3) (-22.3,-15.5) (-22.6,-16.7) (-23.0,-17.9) (-23.3,-19.2) (-23.5,-20.4) (-23.7,-21.6) (-23.9,-22.8) (-24.1,-24.1) (-24.2,-25.3) (-24.3,-26.6) (-24.3,-27.8) (-24.4,-29.1) (-24.3,-29.8)}; \draw[sxc tip] (-21.40,-12.92) -- (-21.57,-13.39);
\draw[sxc fl] plot[smooth] coordinates {(-15.0,0.0) (-14.6,-2.2) (-14.4,-3.4) (-14.1,-4.6) (-13.9,-5.8) (-13.6,-7.1) (-13.2,-8.3) (-12.8,-9.4) (-12.4,-10.6) (-11.9,-11.8) (-11.3,-12.9) (-10.7,-14.0) (-10.0,-15.0) (-9.3,-16.0) (-8.5,-16.9) (-7.5,-17.8) (-6.6,-18.6) (-5.5,-19.3) (-4.4,-19.8) (-3.2,-20.3) (-2.0,-20.6) (-0.8,-20.7) (0.5,-20.8) (1.7,-20.6) (2.9,-20.4) (4.1,-20.0) (5.2,-19.4) (6.3,-18.8) (7.3,-18.0) (8.2,-17.2) (9.1,-16.3) (9.9,-15.3) (10.6,-14.3) (11.2,-13.2) (11.8,-12.1) (12.3,-10.9) (12.7,-9.8) (13.1,-8.6) (13.5,-7.4) (13.8,-6.2) (14.1,-4.9) (14.3,-3.7) (14.6,-2.5) (15.0,0.0)}; \draw[sxc tip] (0.20,-20.85) -- (0.70,-20.77);
\draw[sxc fl] plot[smooth] coordinates {(-15.0,0.0) (-13.3,-1.4) (-12.4,-2.2) (-11.4,-3.0) (-10.4,-3.8) (-9.4,-4.5) (-8.3,-5.2) (-7.3,-5.8) (-6.2,-6.4) (-5.0,-7.0) (-3.9,-7.4) (-2.7,-7.7) (-1.4,-7.9) (-0.2,-8.0) (1.1,-8.0) (2.3,-7.8) (3.5,-7.5) (4.7,-7.1) (5.9,-6.6) (7.0,-6.0) (8.0,-5.4) (9.1,-4.7) (10.1,-4.0) (11.1,-3.2) (12.1,-2.4) (13.0,-1.6) (15.0,0.0)}; \draw[sxc tip] (-0.50,-8.00) -- (0.00,-8.00);
\end{scope}
\draw[bkGray, line width=0.4pt] (-52,-30) rectangle (52,30);
\sxcplus{(-15,0)}{3.2}\sxcminus{(15,0)}{3.2}
\node[sxc lab, anchor=north] at (0,-31) {正と負（同じ大きさ）};
\begin{scope}[shift={(112,0)}]
  \begin{scope}
  \clip (-52,-30) rectangle (52,30);
\draw[sxc eq] plot[smooth, tension=0.6] coordinates {(-11.8,0.0) (-12.0,1.2) (-12.7,2.2) (-13.8,2.9) (-15.0,3.2) (-16.2,2.9) (-17.2,2.2) (-17.9,1.1) (-18.1,-0.1) (-17.9,-1.3) (-17.1,-2.3) (-16.1,-3.0) (-14.9,-3.2) (-13.6,-2.9) (-12.6,-2.1) (-12.0,-1.1) (-11.8,0.0)};
\draw[sxc eq] plot[smooth, tension=0.6] coordinates {(-8.7,0.0) (-8.8,1.2) (-9.2,2.4) (-9.8,3.5) (-10.7,4.4) (-11.7,5.1) (-12.8,5.7) (-14.0,5.9) (-15.3,6.0) (-16.5,5.7) (-17.7,5.3) (-18.7,4.6) (-19.6,3.7) (-20.2,2.6) (-20.6,1.4) (-20.8,0.2) (-20.7,-1.1) (-20.4,-2.3) (-19.8,-3.4) (-19.0,-4.3) (-18.0,-5.1) (-16.9,-5.6) (-15.6,-5.9) (-14.4,-6.0) (-13.2,-5.8) (-12.0,-5.3) (-10.9,-4.7) (-10.1,-3.8) (-9.4,-2.8) (-8.9,-1.6) (-8.7,-0.4) (-8.7,0.0)};
\draw[sxc eq] plot[smooth, tension=0.6] coordinates {(11.8,0.0) (12.0,-1.2) (12.7,-2.2) (13.8,-2.9) (15.0,-3.2) (16.2,-2.9) (17.2,-2.2) (17.9,-1.1) (18.1,0.1) (17.9,1.3) (17.1,2.3) (16.1,3.0) (14.9,3.2) (13.6,2.9) (12.6,2.1) (12.0,1.1) (11.8,0.0)};
\draw[sxc eq] plot[smooth, tension=0.6] coordinates {(8.7,0.0) (8.8,-1.2) (9.2,-2.4) (9.8,-3.5) (10.7,-4.4) (11.7,-5.1) (12.8,-5.7) (14.0,-5.9) (15.3,-6.0) (16.5,-5.7) (17.7,-5.3) (18.7,-4.6) (19.6,-3.7) (20.2,-2.6) (20.6,-1.4) (20.8,-0.2) (20.7,1.1) (20.4,2.3) (19.8,3.4) (19.0,4.3) (18.0,5.1) (16.9,5.6) (15.6,5.9) (14.4,6.0) (13.2,5.8) (12.0,5.3) (10.9,4.7) (10.1,3.8) (9.4,2.8) (8.9,1.6) (8.7,0.4) (8.7,0.0)};
\draw[sxc eq] plot[smooth, tension=0.6] coordinates {(28.0,0.0) (28.0,1.2) (27.8,2.5) (27.5,3.7) (27.1,4.9) (26.6,6.0) (26.1,7.1) (25.4,8.2) (24.6,9.2) (23.7,10.1) (22.8,10.9) (21.8,11.7) (20.7,12.3) (19.6,12.9) (18.5,13.4) (17.3,13.8) (16.1,14.1) (14.9,14.3) (13.6,14.5) (12.4,14.5) (11.1,14.5) (9.9,14.4) (8.6,14.3) (7.4,14.1) (6.2,14.0) (4.9,13.7) (3.7,13.5) (2.5,13.4) (1.2,13.3) (-0.0,13.2) (-1.3,13.3) (-2.5,13.4) (-3.8,13.6) (-5.0,13.8) (-6.2,14.0) (-7.5,14.2) (-8.7,14.3) (-10.0,14.5) (-11.2,14.5) (-12.5,14.5) (-13.7,14.5) (-14.9,14.3) (-16.2,14.1) (-17.4,13.8) (-18.6,13.4) (-19.7,12.9) (-20.8,12.3) (-21.9,11.6) (-22.9,10.9) (-23.8,10.0) (-24.6,9.1) (-25.4,8.1) (-26.1,7.1) (-26.7,6.0) (-27.2,4.8) (-27.5,3.6) (-27.8,2.4) (-28.0,1.2) (-28.0,-0.1) (-28.0,-1.3) (-27.8,-2.6) (-27.5,-3.8) (-27.1,-5.0) (-26.6,-6.1) (-26.0,-7.2) (-25.3,-8.3) (-24.5,-9.2) (-23.7,-10.1) (-22.7,-11.0) (-21.7,-11.7) (-20.7,-12.4) (-19.6,-12.9) (-18.4,-13.4) (-17.2,-13.8) (-16.0,-14.1) (-14.8,-14.3) (-13.5,-14.5) (-12.3,-14.5) (-11.0,-14.5) (-9.8,-14.4) (-8.5,-14.3) (-7.3,-14.1) (-6.1,-13.9) (-4.8,-13.7) (-3.6,-13.5) (-2.4,-13.4) (-1.1,-13.3) (0.1,-13.2) (1.4,-13.3) (2.6,-13.4) (3.9,-13.6) (5.1,-13.8) (6.3,-14.0) (7.6,-14.2) (8.8,-14.3) (10.0,-14.5) (11.3,-14.5) (12.5,-14.5) (13.8,-14.5) (15.0,-14.3) (16.3,-14.1) (17.5,-13.7) (18.6,-13.3) (19.8,-12.8) (20.9,-12.2) (21.9,-11.6) (22.9,-10.8) (23.9,-10.0) (24.7,-9.0) (25.5,-8.0) (26.1,-7.0) (26.7,-5.9) (27.2,-4.7) (27.6,-3.5) (27.8,-2.3) (28.0,-1.1) (28.0,0.0)};
\draw[sxc eq] plot[smooth, tension=0.6] coordinates {(35.0,0.0) (35.0,1.2) (34.9,2.5) (34.7,3.7) (34.4,5.0) (34.1,6.2) (33.7,7.4) (33.3,8.5) (32.8,9.7) (32.2,10.8) (31.5,11.8) (30.8,12.9) (30.1,13.9) (29.3,14.8) (28.4,15.7) (27.5,16.6) (26.6,17.4) (25.6,18.2) (24.6,18.9) (23.5,19.6) (22.4,20.2) (21.3,20.8) (20.2,21.3) (19.0,21.7) (17.8,22.2) (16.6,22.5) (15.4,22.9) (14.2,23.1) (13.0,23.4) (11.8,23.6) (10.5,23.8) (9.3,23.9) (8.0,24.0) (6.8,24.1) (5.5,24.2) (4.3,24.2) (3.0,24.3) (1.8,24.3) (0.5,24.3) (-0.7,24.3) (-2.0,24.3) (-3.2,24.3) (-4.5,24.2) (-5.7,24.2) (-7.0,24.1) (-8.2,24.0) (-9.4,23.9) (-10.7,23.8) (-11.9,23.6) (-13.2,23.4) (-14.4,23.1) (-15.6,22.8) (-16.8,22.5) (-18.0,22.1) (-19.2,21.7) (-20.3,21.2) (-21.5,20.7) (-22.6,20.1) (-23.6,19.5) (-24.7,18.8) (-25.7,18.1) (-26.7,17.3) (-27.6,16.5) (-28.5,15.6) (-29.4,14.7) (-30.2,13.7) (-30.9,12.7) (-31.6,11.7) (-32.3,10.6) (-32.8,9.5) (-33.3,8.4) (-33.8,7.2) (-34.2,6.0) (-34.5,4.8) (-34.7,3.6) (-34.9,2.3) (-35.0,1.1) (-35.0,-0.2) (-35.0,-1.4) (-34.8,-2.7) (-34.6,-3.9) (-34.4,-5.1) (-34.1,-6.3) (-33.7,-7.5) (-33.2,-8.7) (-32.7,-9.8) (-32.1,-10.9) (-31.4,-12.0) (-30.7,-13.0) (-30.0,-14.0) (-29.2,-14.9) (-28.3,-15.9) (-27.4,-16.7) (-26.4,-17.5) (-25.5,-18.3) (-24.4,-19.0) (-23.4,-19.7) (-22.3,-20.3) (-21.2,-20.8) (-20.0,-21.3) (-18.9,-21.8) (-17.7,-22.2) (-16.5,-22.6) (-15.3,-22.9) (-14.1,-23.2) (-12.8,-23.4) (-11.6,-23.6) (-10.4,-23.8) (-9.1,-23.9) (-7.9,-24.0) (-6.6,-24.1) (-5.4,-24.2) (-4.1,-24.3) (-2.9,-24.3) (-1.6,-24.3) (-0.4,-24.3) (0.9,-24.3) (2.1,-24.3) (3.4,-24.3) (4.6,-24.2) (5.9,-24.2) (7.1,-24.1) (8.4,-24.0) (9.6,-23.9) (10.8,-23.7) (12.1,-23.6) (13.3,-23.3) (14.5,-23.1) (15.8,-22.8) (17.0,-22.4) (18.1,-22.1) (19.3,-21.6) (20.5,-21.1) (21.6,-20.6) (22.7,-20.0) (23.8,-19.4) (24.8,-18.7) (25.9,-18.0) (26.8,-17.2) (27.8,-16.4) (28.7,-15.5) (29.5,-14.6) (30.3,-13.6) (31.0,-12.6) (31.7,-11.6) (32.3,-10.5) (32.9,-9.4) (33.4,-8.2) (33.8,-7.0) (34.2,-5.8) (34.5,-4.6) (34.7,-3.4) (34.9,-2.2) (35.0,-0.9) (35.0,0.0)};
\draw[sxc eq] plot[smooth, tension=0.6] coordinates {(27.3,-30.1) (28.4,-29.4) (29.4,-28.7) (30.4,-28.0) (31.4,-27.2) (32.4,-26.5) (33.3,-25.6) (34.3,-24.8) (35.1,-23.9) (36.0,-23.0) (36.8,-22.0) (37.6,-21.1) (38.4,-20.1) (39.1,-19.1) (39.8,-18.0) (40.4,-16.9) (41.0,-15.8) (41.6,-14.7) (42.1,-13.6) (42.6,-12.4) (43.0,-11.3) (43.4,-10.1) (43.8,-8.9) (44.1,-7.7) (44.4,-6.5) (44.6,-5.2) (44.8,-4.0) (44.9,-2.7) (45.0,-1.5) (45.0,-0.2) (45.0,1.0) (44.9,2.2) (44.8,3.5) (44.7,4.7) (44.5,6.0) (44.2,7.2) (43.9,8.4) (43.6,9.6) (43.2,10.8) (42.8,12.0) (42.3,13.1) (41.8,14.3) (41.3,15.4) (40.7,16.5) (40.0,17.6) (39.4,18.6) (38.7,19.7) (37.9,20.7) (37.1,21.7) (36.3,22.6) (35.5,23.5) (34.6,24.4) (33.7,25.3) (32.8,26.1) (31.8,26.9) (30.8,27.7) (29.8,28.4) (28.8,29.2) (27.8,29.8) (27.3,30.1)};
\draw[sxc fl] plot[smooth] coordinates {(-15.0,0.0) (-13.3,1.4) (-12.4,2.2) (-11.4,3.1) (-10.5,3.9) (-9.6,4.8) (-8.8,5.7) (-8.0,6.6) (-7.2,7.7) (-6.6,8.7) (-6.0,9.9) (-5.6,11.0) (-5.3,12.2) (-5.0,13.4) (-4.8,14.7) (-4.7,15.9) (-4.6,17.2) (-4.5,18.4) (-4.5,19.7) (-4.5,20.9) (-4.5,22.2) (-4.5,23.4) (-4.6,24.7) (-4.6,25.9) (-4.7,27.2) (-4.8,28.4) (-4.9,29.7) (-4.9,29.9)}; \draw[sxc tip] (-5.37,11.91) -- (-5.25,12.39);
\draw[sxc fl] plot[smooth] coordinates {(-15.0,0.0) (-14.6,2.2) (-14.4,3.4) (-14.2,4.6) (-14.1,5.9) (-13.9,7.1) (-13.8,8.4) (-13.8,9.6) (-13.8,10.9) (-13.8,12.1) (-13.9,13.4) (-14.0,14.6) (-14.1,15.8) (-14.3,17.1) (-14.5,18.3) (-14.8,19.5) (-15.0,20.8) (-15.3,22.0) (-15.6,23.2) (-15.9,24.4) (-16.3,25.6) (-16.6,26.8) (-17.0,28.0) (-17.4,29.2) (-17.6,29.9)}; \draw[sxc tip] (-13.98,14.30) -- (-14.02,14.80);
\draw[sxc fl] plot[smooth] coordinates {(-15.0,0.0) (-16.1,1.9) (-16.7,3.0) (-17.4,4.1) (-18.0,5.1) (-18.7,6.2) (-19.4,7.2) (-20.1,8.3) (-20.8,9.3) (-21.5,10.3) (-22.3,11.3) (-23.0,12.3) (-23.8,13.3) (-24.6,14.3) (-25.4,15.2) (-26.2,16.2) (-27.0,17.1) (-27.8,18.1) (-28.6,19.0) (-29.4,20.0) (-30.3,20.9) (-31.1,21.8) (-32.0,22.7) (-32.8,23.7) (-33.7,24.6) (-34.5,25.5) (-35.4,26.4) (-36.3,27.3) (-37.1,28.2) (-38.0,29.1) (-38.9,29.9)}; \draw[sxc tip] (-24.40,14.08) -- (-24.73,14.45);
\draw[sxc fl] plot[smooth] coordinates {(-15.0,0.0) (-17.1,0.8) (-18.2,1.2) (-19.4,1.6) (-20.6,2.0) (-21.8,2.4) (-23.0,2.8) (-24.1,3.2) (-25.3,3.6) (-26.5,4.0) (-27.7,4.4) (-28.9,4.8) (-30.1,5.2) (-31.3,5.5) (-32.5,5.9) (-33.7,6.3) (-34.9,6.7) (-36.1,7.0) (-37.3,7.4) (-38.5,7.7) (-39.7,8.1) (-40.9,8.4) (-42.1,8.8) (-43.3,9.1) (-44.5,9.5) (-45.7,9.8) (-46.9,10.2) (-48.1,10.5) (-49.3,10.8) (-50.5,11.2) (-51.7,11.5) (-51.9,11.6)}; \draw[sxc tip] (-32.22,5.81) -- (-32.69,5.96);
\draw[sxc fl] plot[smooth] coordinates {(-15.0,0.0) (-17.1,-0.8) (-18.2,-1.2) (-19.4,-1.6) (-20.6,-2.0) (-21.8,-2.4) (-23.0,-2.8) (-24.1,-3.2) (-25.3,-3.6) (-26.5,-4.0) (-27.7,-4.4) (-28.9,-4.8) (-30.1,-5.2) (-31.3,-5.5) (-32.5,-5.9) (-33.7,-6.3) (-34.9,-6.7) (-36.1,-7.0) (-37.3,-7.4) (-38.5,-7.7) (-39.7,-8.1) (-40.9,-8.4) (-42.1,-8.8) (-43.3,-9.1) (-44.5,-9.5) (-45.7,-9.8) (-46.9,-10.2) (-48.1,-10.5) (-49.3,-10.8) (-50.5,-11.2) (-51.7,-11.5) (-51.9,-11.6)}; \draw[sxc tip] (-32.22,-5.81) -- (-32.69,-5.96);
\draw[sxc fl] plot[smooth] coordinates {(-15.0,0.0) (-16.1,-1.9) (-16.7,-3.0) (-17.4,-4.1) (-18.0,-5.1) (-18.7,-6.2) (-19.4,-7.2) (-20.1,-8.3) (-20.8,-9.3) (-21.5,-10.3) (-22.3,-11.3) (-23.0,-12.3) (-23.8,-13.3) (-24.6,-14.3) (-25.4,-15.2) (-26.2,-16.2) (-27.0,-17.1) (-27.8,-18.1) (-28.6,-19.0) (-29.4,-20.0) (-30.3,-20.9) (-31.1,-21.8) (-32.0,-22.7) (-32.8,-23.7) (-33.7,-24.6) (-34.5,-25.5) (-35.4,-26.4) (-36.3,-27.3) (-37.1,-28.2) (-38.0,-29.1) (-38.9,-29.9)}; \draw[sxc tip] (-24.40,-14.08) -- (-24.73,-14.45);
\draw[sxc fl] plot[smooth] coordinates {(-15.0,0.0) (-14.6,-2.2) (-14.4,-3.4) (-14.2,-4.6) (-14.1,-5.9) (-13.9,-7.1) (-13.8,-8.4) (-13.8,-9.6) (-13.8,-10.9) (-13.8,-12.1) (-13.9,-13.4) (-14.0,-14.6) (-14.1,-15.8) (-14.3,-17.1) (-14.5,-18.3) (-14.8,-19.5) (-15.0,-20.8) (-15.3,-22.0) (-15.6,-23.2) (-15.9,-24.4) (-16.3,-25.6) (-16.6,-26.8) (-17.0,-28.0) (-17.4,-29.2) (-17.6,-29.9)}; \draw[sxc tip] (-13.98,-14.30) -- (-14.02,-14.80);
\draw[sxc fl] plot[smooth] coordinates {(-15.0,0.0) (-13.3,-1.4) (-12.4,-2.2) (-11.4,-3.1) (-10.5,-3.9) (-9.6,-4.8) (-8.8,-5.7) (-8.0,-6.6) (-7.2,-7.7) (-6.6,-8.7) (-6.0,-9.9) (-5.6,-11.0) (-5.3,-12.2) (-5.0,-13.4) (-4.8,-14.7) (-4.7,-15.9) (-4.6,-17.2) (-4.5,-18.4) (-4.5,-19.7) (-4.5,-20.9) (-4.5,-22.2) (-4.5,-23.4) (-4.6,-24.7) (-4.6,-25.9) (-4.7,-27.2) (-4.8,-28.4) (-4.9,-29.7) (-4.9,-29.9)}; \draw[sxc tip] (-5.37,-11.91) -- (-5.25,-12.39);
\draw[sxc fl] plot[smooth] coordinates {(15.0,0.0) (13.3,1.4) (12.4,2.2) (11.4,3.1) (10.5,3.9) (9.6,4.8) (8.8,5.7) (8.0,6.6) (7.2,7.7) (6.6,8.7) (6.0,9.9) (5.6,11.0) (5.3,12.2) (5.0,13.4) (4.8,14.7) (4.7,15.9) (4.6,17.2) (4.5,18.4) (4.5,19.7) (4.5,20.9) (4.5,22.2) (4.5,23.4) (4.6,24.7) (4.6,25.9) (4.7,27.2) (4.8,28.4) (4.9,29.7) (4.9,29.9)}; \draw[sxc tip] (5.37,11.91) -- (5.25,12.39);
\draw[sxc fl] plot[smooth] coordinates {(15.0,0.0) (14.6,2.2) (14.4,3.4) (14.2,4.6) (14.1,5.9) (13.9,7.1) (13.8,8.4) (13.8,9.6) (13.8,10.9) (13.8,12.1) (13.9,13.4) (14.0,14.6) (14.1,15.8) (14.3,17.1) (14.5,18.3) (14.8,19.5) (15.0,20.8) (15.3,22.0) (15.6,23.2) (15.9,24.4) (16.3,25.6) (16.6,26.8) (17.0,28.0) (17.4,29.2) (17.6,29.9)}; \draw[sxc tip] (13.98,14.30) -- (14.02,14.80);
\draw[sxc fl] plot[smooth] coordinates {(15.0,0.0) (16.1,1.9) (16.7,3.0) (17.4,4.1) (18.0,5.1) (18.7,6.2) (19.4,7.2) (20.1,8.3) (20.8,9.3) (21.5,10.3) (22.3,11.3) (23.0,12.3) (23.8,13.3) (24.6,14.3) (25.4,15.2) (26.2,16.2) (27.0,17.1) (27.8,18.1) (28.6,19.0) (29.4,20.0) (30.3,20.9) (31.1,21.8) (32.0,22.7) (32.8,23.7) (33.7,24.6) (34.5,25.5) (35.4,26.4) (36.3,27.3) (37.1,28.2) (38.0,29.1) (38.9,29.9)}; \draw[sxc tip] (24.40,14.08) -- (24.73,14.45);
\draw[sxc fl] plot[smooth] coordinates {(15.0,0.0) (16.1,-1.9) (16.7,-3.0) (17.4,-4.1) (18.0,-5.1) (18.7,-6.2) (19.4,-7.2) (20.1,-8.3) (20.8,-9.3) (21.5,-10.3) (22.3,-11.3) (23.0,-12.3) (23.8,-13.3) (24.6,-14.3) (25.4,-15.2) (26.2,-16.2) (27.0,-17.1) (27.8,-18.1) (28.6,-19.0) (29.4,-20.0) (30.3,-20.9) (31.1,-21.8) (32.0,-22.7) (32.8,-23.7) (33.7,-24.6) (34.5,-25.5) (35.4,-26.4) (36.3,-27.3) (37.1,-28.2) (38.0,-29.1) (38.9,-29.9)}; \draw[sxc tip] (24.40,-14.08) -- (24.73,-14.45);
\draw[sxc fl] plot[smooth] coordinates {(15.0,0.0) (14.6,-2.2) (14.4,-3.4) (14.2,-4.6) (14.1,-5.9) (13.9,-7.1) (13.8,-8.4) (13.8,-9.6) (13.8,-10.9) (13.8,-12.1) (13.9,-13.4) (14.0,-14.6) (14.1,-15.8) (14.3,-17.1) (14.5,-18.3) (14.8,-19.5) (15.0,-20.8) (15.3,-22.0) (15.6,-23.2) (15.9,-24.4) (16.3,-25.6) (16.6,-26.8) (17.0,-28.0) (17.4,-29.2) (17.6,-29.9)}; \draw[sxc tip] (13.98,-14.30) -- (14.02,-14.80);
\draw[sxc fl] plot[smooth] coordinates {(15.0,0.0) (13.3,-1.4) (12.4,-2.2) (11.4,-3.1) (10.5,-3.9) (9.6,-4.8) (8.8,-5.7) (8.0,-6.6) (7.2,-7.7) (6.6,-8.7) (6.0,-9.9) (5.6,-11.0) (5.3,-12.2) (5.0,-13.4) (4.8,-14.7) (4.7,-15.9) (4.6,-17.2) (4.5,-18.4) (4.5,-19.7) (4.5,-20.9) (4.5,-22.2) (4.5,-23.4) (4.6,-24.7) (4.6,-25.9) (4.7,-27.2) (4.8,-28.4) (4.9,-29.7) (4.9,-29.9)}; \draw[sxc tip] (5.37,-11.91) -- (5.25,-12.39);
  \end{scope}
  \draw[bkGray, line width=0.4pt] (-52,-30) rectangle (52,30);
  \sxcplus{(-15,0)}{3.2}\sxcplus{(15,0)}{3.2}
  \node[sxc lab, anchor=north] at (0,-31) {正と正（同じ大きさ）};
\end{scope}
\end{tikzpicture}
\bkfigcaption{2つの点電荷の電気力線（実線。矢じりは電場の向き）と等電位線（破線）。どこでも、電気力線と等電位線は垂直に交わる。左の、2つのまん中の縦の破線は、電位が0の等電位線。Python で計算してかいた。}\label{ph:電-2:fig-toudeni}
\end{bkfigure}

\phboxhead{電気力線と等電位面は垂直}
等電位面にそって電荷を動かしても、電位が変わらないので、電場の力は仕事をしません。力が仕事をしないのは、力が動く向きに垂直なときです（力-4）。だから、\textbf{電場（電気力線）は、等電位面にいつも垂直}です。\label{ph:電-2:suichoku}

\phboxhead{電場は、電位が急に下がる向きに強い}
電気力線にそって距離 $\Delta s$ 進むと、電位が $\Delta\phi$ だけ変わるとします（下がるとき $\Delta\phi<0$）。そのあいだに電場がする仕事 $qE\,\Delta s$ は、位置エネルギーの減り $-q\,\Delta\phi$ に等しいので、
\pheq{E=-\frac{\Delta\phi}{\Delta s}}{V/m ＝ V ÷ m}
です。電場は、電位が下がる向きを向き、\textbf{下がり方が急な所ほど強い}のです。地図の上で、等高線がこみ合っている所ほど坂が急で、ボールがよく転がるのと同じです。図\ref{ph:電-2:fig-toudeni}でも、電荷の近くの等電位線がこみ合っている所で、電場が強くなっています。一様な電場の $E=V/d$（⑥）は、下がり方が一定のときの、この式です。

\begin{phmijika}[鳥が電線にとまっても大丈夫なわけ（入口）]
電場の中に金属を置くと、金属の中の自由電子が動いて、金属全体が1つの等電位になります（電-3）。電線にとまった鳥の両足は、同じ電線の上の、ほとんど同じ電位の2点にあるので、足のあいだの電圧はほぼ0です。だから、鳥の体にはほとんど電流が流れません（くわしくは電-5 のはてな？）。ただし、2本の電線や、電線と電柱に同時に触れると、大きな電圧がかかるので、とても危険です。
\end{phmijika}

\begin{phmodel}[止まった電荷の電場と電位]
\phignore{電荷が動くことによる磁気の効果。空気が電場をわずかに弱める効果（空気の比誘電率は約1.0005：電-3）。}
\phassume{電荷は止まっている（静電場）。電荷の周りは真空。}
\phlimit{電荷が速く動いたり、電流が急に変わったりするとき（電場が保存力でなくなる：電-8）。水やガラスなどの物の中（電-3）。}
\end{phmodel}

\phkokomade

\Needspace{10\baselineskip}\phstage{shiki}{ベクトル解析の式で書く}{5}

\bksubsection{8}{ガウスの法則の積分形と微分形}
\label{ph:電-2:sec8}
③の「袋から出ていく電気力線の本数」は、数学の言葉では、閉じた面の上の\textbf{流束}（8-8）です。袋の面 $S$ を小さなタイルに分け、タイルの面積 $\dd S$、外向きの長さ1の矢印（法線）を $\vec n$ とすると、タイルを通る本数は $\vec E\cdot\vec n\,\dd S$ です（電場の、面に垂直な成分×面積）。これを袋全体でたしたものが $\Phi_E$ なので、ガウスの法則は
\pheq{\oiint_S\vec E\cdot\vec n\,\dd S=\frac{Q_{\text{中}}}{\varepsilon_0}}{N·m²/C}
\label{ph:電-2:sekibun}と書けます（\textbf{積分形}。8-11⑦）。④で使った「垂直で同じ強さなら $E\times$面積、平行なら0」は、この積分が簡単になる場合でした（8-8の場合 D）。

\phboxhead{電荷が空間に広がっているとき}
電荷が、\qty{1}{m^3}あたり $\rho$〔C/m$^3$〕の密度で空間に広がっているとき、袋の中の電荷は体積の積分 $Q_{\text{中}}=\iiint_{\text{中}}\rho\,\dd V$ です（8-6。$\dd V$ は小さな体積で、電圧の $V$ とは別）。いっぽう、ガウスの定理（8-11⑤）で、左辺の面の積分は、発散の体積の積分に取りかえられます。
\[\oiint_S\vec E\cdot\vec n\,\dd S=\iiint_{\text{中}}\mathrm{div}\,\vec E\,\dd V=\iiint_{\text{中}}\frac{\rho}{\varepsilon_0}\,\dd V\]
これが、どんな小さな袋でも成り立つので、積分の中どうしが等しく、
\pheq{\mathrm{div}\,\vec E=\frac{\rho}{\varepsilon_0}}{$\dfrac{\text{N/C}}{\text{m}}$ ＝ $\dfrac{\text{C/m}^3}{\text{C}^2/(\text{N·m}^2)}$}
\label{ph:電-2:bibun}です（\textbf{微分形}。8-9⑬）。\bkidx{はっさん}{発散（電場）}「電場のわき出しの強さ（発散）は、その場所の電荷の密度に比例する」と読めます。電荷のない所では $\mathrm{div}\,\vec E=0$ で、電気力線は生まれも消えもしません。

\begin{phyomi}[$\mathrm{div}\,\vec E=\rho/\varepsilon_0$]
積分形は「袋全体で数える」、微分形は「1つの点の周りの小さな箱で数える」形です。同じ法則を、大きい袋で見るか、小さい箱で見るかの違いです。対称な問題は積分形、コンピュータで場所ごとに計算するときは微分形が便利です。
\end{phyomi}

\noindent 一様に帯電した球の中の電場 $E=kQr/R^3$（④）で、確かめましょう。電場が中心からの距離 $r$ だけで決まり、放射状に向いているとき、発散は $\dfrac{1}{r^2}\dfrac{\dd}{\dd r}(r^2E)$ です（8-12⑦ の球の式と同じ考え）。
\[\mathrm{div}\,\vec E=\frac{1}{r^2}\,\frac{\dd}{\dd r}\Bigl(\frac{kQ}{R^3}r^3\Bigr)=\frac{3kQ}{R^3}=\frac{3Q}{4\pi\varepsilon_0R^3}=\frac{1}{\varepsilon_0}\times\frac{Q}{\frac43\pi R^3}\]
最後の $Q/(\tfrac43\pi R^3)$ は、球の中の電荷の密度 $\rho$ です。微分形が成り立っています。球の外では $r^2E=kQ$ が一定なので、発散は0です。

%% ---------------------------------------------------------------
\bksubsection{9}{電場は保存力の場：$\mathrm{rot}\,\vec E=\vec 0$ と $\vec E=-\mathrm{grad}\,\phi$}
\label{ph:電-2:sec9}
⑤では、電位を「無限に遠い所から運んでくる仕事」で決めました。数学の言葉で書くと、点 P の電位は、電場の\textbf{線積分}（8-7）です。
\pheq{\phi(\mathrm P)=-\int_{\infty}^{\mathrm P}\vec E\cdot\dd\vec r}{V ＝ V/m × m}
この値が、途中の道によらないことが大切です。止まった電荷の電場は、点電荷の電場の重ね合わせで、点電荷の電場は中心向きの力（$1/r^2$）なので、保存力です（力-4、8-7⑨）。道によらないことは、「閉じた道を1周すると仕事が0」と同じです。
\[\oint_C\vec E\cdot\dd\vec r=0\]
ストークスの定理（8-11⑨）で、これがどんな小さな輪でも成り立つことは、
\pheq{\mathrm{rot}\,\vec E=\vec 0}{1/m × V/m}
と同じです。そして、電位から電場を作る式は、⑦の $E=-\Delta\phi/\Delta s$ を3次元の矢印にした
\pheq{\vec E=-\mathrm{grad}\,\phi=-\Bigl(\frac{\partial\phi}{\partial x},\ \frac{\partial\phi}{\partial y},\ \frac{\partial\phi}{\partial z}\Bigr)}{V/m ＝ V ÷ m}
\label{ph:電-2:grad}です（8-4⑧）。勾配は等値面に垂直なので（8-4②）、高校の「等電位面と電気力線は垂直」は、この式の言いかえです。

\begin{phflow}[電位と電場]
\phflowrow{$\vec E=-\mathrm{grad}\,\phi$}{一様な電場\\$\phi$ が $x$ の一次式}{$V=Ed$}
\phflowrow{$\phi=-\displaystyle\int_\infty^r\vec E\cdot\dd\vec r$}{点電荷\\$E=kQ/r^2$}{$\phi=k\,\dfrac{Q}{r}$}
\end{phflow}

\begin{seikaku}
$\mathrm{rot}\,\vec E=\vec0$ は、電荷が止まっているとき（静電場）だけの式です。磁場が時間とともに変わると、電場はうずを巻き（$\mathrm{rot}\,\vec E\neq\vec0$）、1周の仕事が0でなくなります。これが発電機のしくみ（電磁誘導：電-8）です。回路の「1周すると電圧の和は0」（キルヒホッフの電圧則：電-4）は、$\oint\vec E\cdot\dd\vec r=0$ から出てきます。
\end{seikaku}

%% ---------------------------------------------------------------
\bksubsection{10}{ポアソン方程式とラプラス方程式}
\label{ph:電-2:sec10}
⑧の $\mathrm{div}\,\vec E=\rho/\varepsilon_0$ に、⑨の $\vec E=-\mathrm{grad}\,\phi$ を入れます。$\mathrm{div}\,\mathrm{grad}=\nabla^2$（ラプラシアン：8-12⑤）なので、
\pheq{\nabla^2\phi=-\frac{\rho}{\varepsilon_0}}{V/m² ＝ (C/m³)÷(F/m)}
です。これを\textbf{ポアソン方程式}といいます。\bkidx{ぽあそんほうていしき}{ポアソン方程式}電荷のない所（$\rho=0$）では、
\pheq{\nabla^2\phi=0}{V/m²}
\label{ph:電-2:laplace}で、これが\textbf{ラプラス方程式}です（8-12⑥）。\bkidx{らぷらすほうていしき}{ラプラス方程式}8-12①で学んだように、$\nabla^2\phi=0$ は「\textbf{どの点の電位も、周りの点の電位の平均と同じ}」という意味です。電荷のない空間では、電位は、周りにならされた、なめらかな形になります。

\phboxhead{電位の山の頂上は、電荷のない所にはない}
どの点も周りの平均と同じなら、周りより高い点（山の頂上）も、低い点（谷の底）もありません（8-12⑥の最大値の原理）。電位のいちばん高い所と低い所は、電荷のある所か、境目（電極）にしかないのです。このことから、おもしろいことがわかります。正の電荷を、止まった電荷がつくる電場だけで、空中の1点にじっと止めておくことはできません。止めておくには、その点が電位の谷の底になっていなければなりませんが、電荷のない所に谷の底はないからです（専門の本では「アーンショーの定理」とよばれます）。

\begin{motto}[5]
\phboxhead{一様に帯電した球の電位を、ポアソン方程式で求める}
④の一様な球（電荷 $Q$、半径 $R$）の中では、電荷の密度が $\rho=Q/(\tfrac43\pi R^3)$ で一定です。電位は中心からの距離 $r$ だけで決まるので、8-12⑦の球の式が使えます。
\[\frac{1}{r^2}\,\frac{\dd}{\dd r}\Bigl(r^2\,\frac{\dd\phi}{\dd r}\Bigr)=-\frac{\rho}{\varepsilon_0}\]
両辺に $r^2$ をかけて積分すると $r^2\,\dd\phi/\dd r=-\rho r^3/(3\varepsilon_0)$（中心で電場が無限大にならないように、積分の定数は0）。もう一度積分して、表面で $\phi=kQ/R$（外の電位とつながる）とすると、
\[\phi=\frac{kQ}{2R^3}\,(3R^2-r^2)\qquad(r\leqq R)\]
です。$-\dd\phi/\dd r=kQr/R^3$ で、④の電場と合います。中心の電位は表面の1.5倍です。$Q=\qty{1e-8}{C}$、$R=\qty{0.1}{m}$ なら、表面\qty{899}{V}、中心は約\qty{1350}{V}です。電位の山の頂上は、電荷のある球の中（中心）にあり、⑩の「電荷のない所に頂上はない」と矛盾しません。
\end{motto}

\phboxhead{答えは1つ（一意性）}
電荷の置き方と、周りの電極（境目）の電位が決まれば、ポアソン方程式の答え（電位の分布）は\textbf{ただ1つ}に決まります。これを\textbf{解の一意性}といいます。\bkidx{かいのいちいせい}{解の一意性}だから、どんな方法でも、条件をみたす答えを1つ見つければ、それが正しい答えです（電-3 ⑫の鏡像法は、このことを使う）。

\begin{seikaku}
一意性のわけ：同じ条件をみたす答えが2つ $\phi_1$・$\phi_2$ あったとします。差 $\psi=\phi_1-\phi_2$ は、電荷のない式 $\nabla^2\psi=0$ をみたし（ポアソン方程式の右辺どうしが消える）、境目では0です（どちらも同じ電位だから）。最大値の原理で、$\psi$ のいちばん大きい値も小さい値も境目の0なので、$\psi$ はどこでも0。つまり $\phi_1=\phi_2$ です。
\end{seikaku}

\noindent 一意性から、例題12の大切な結果が出ます。中に電荷のない、閉じた金属の箱を考えます。金属の箱は全体が1つの電位 $\phi_0$ になります（電-3 ③）。「箱の中でどこでも $\phi=\phi_0$」は、ラプラス方程式（$\nabla^2\phi_0=0$）も境目の条件もみたすので、一意性から、これが答えです。電位が一定なので、箱の中の電場は0です。外にどんな電荷があっても変わりません（電-3 のはてな？、静電遮へい）。

%% ---------------------------------------------------------------
\Needspace{12\baselineskip}\bksubsection{11}{表計算でラプラス方程式を解く}
\label{ph:電-2:sec11}
ラプラス方程式は、ふつう式では解けません。コンピュータで、空間をマスに分け、「どの点も周り4点の平均」になるまで、ならすことを繰り返します。8-12③で、板の温度を同じ方法で求めました。ここでは、温度を電位に読みかえます。

正方形の金属の箱（たての断面）を考えます。上の辺だけ、ほかの辺と切りはなした金属の板で、\qty{100}{V}にします。ほかの3つの辺は\qty{0}{V}です。断面を $4\times4$ のマスに分け、マスの角の $5\times5$ 個の点のうち、まん中の $3\times3$ の点の電位を求めます。
\begin{enumerate}[label={\bkcircled[bkGreen]{\arabic*}}, leftmargin=2.4\zw, itemsep=1pt, topsep=2pt]
\item 表計算の1〜5行めに、はじめの表を入れる（上のふち100、ほかのふち0、まん中の9点は0）。
\item 7〜11行めのまん中の9マスに、1つ前の表の「北・西・東・南」の平均の式を入れる（B8 に \texttt{=AVERAGE(B1,A2,C2,B3)}）。ふちは1つ前の表と同じ値。
\item 6行ずつ下へコピーして、ならす回数をふやす。表が変わらなくなったら、それが答え。
\end{enumerate}

\begin{bkcenter}
\bksheet{5}{14}{%
1 & 0 & 100 & 100 & 100 & 0 \\
2 & 0 & 0 & 0 & 0 & 0 \\
\SetCell[c=6]{c} 3〜5行めも0 & & & & & \\
7 & 0 & 100 & 100 & 100 & 0 \\
8 & 0 & 25 & 25 & 25 & 0 \\
\SetCell[c=6]{c} 9〜11行めも0 & & & & & \\}
\bktabcaption{はじめの表（1〜5行め）と、1回ならした表（7〜11行め。6行めは空けておく）。上の段は、北どなりが\qty{100}{V}なので $(100+0+0+0)\div4=25$。}\label{ph:電-2:tab-sheet}
\end{bkcenter}

\noindent ならすことを50回ほど繰り返すと、表は変わらなくなり、次の電位〔V〕になります（8-12③の表と同じ数）。
\begin{bkcenter}
\begin{bktable}{colspec={Q[c,m,wd=26mm] *{3}{Q[c,m,wd=20mm]}}, rowsep=1.5pt, row{1}={bg=bkTheme!10, font=\sffamily\small}}
 & 左 & まん中 & 右 \\
上の段 & 42.9 & 52.7 & 42.9 \\
まん中の段 & 18.75 & 25.0 & 18.75 \\
下の段 & 7.14 & 9.82 & 7.14 \\
\end{bktable}
\bktabcaption{落ち着いた電位（上の辺\qty{100}{V}、ほかは\qty{0}{V}）。どの点も、周り4点の平均と同じ。}\label{ph:電-2:tab-kotae}
\end{bkcenter}

\noindent 同じ繰り返しを Python で書くと、次のようになります（8-12③ と同じ形。\texttt{V[1:4, 1:4]} は、まん中の $3\times3$ の点）。
\begin{phcode}[Python]
import numpy as np
V = np.zeros((5, 5))            \# 5×5 の点の電位。はじめは全部0
V[0, 1:4] = 100                 \# 上のふち（角はのぞく）を 100 V
for n in range(50):             \# 50回ならす
    W = V.copy()
    W[1:4, 1:4] = (V[0:3, 1:4] + V[2:5, 1:4]
                   + V[1:4, 0:3] + V[1:4, 2:5]) / 4
    V = W
print(V[1:4, 1:4].round(2))     \# まん中の9点（小数第2位まで）
\end{phcode}
\begin{phcodeout}
[[42.86 52.68 42.86]
 [18.75 25.   18.75]
 [ 7.14  9.82  7.14]]
\end{phcodeout}
4つの「となり」（北・南・西・東）の配列をたして4でわる1行が、表計算の \texttt{=AVERAGE(\ldots)} を9マスに貼ったのと同じ働きをします。\texttt{np.zeros((5, 5))} の5を大きくすれば、マスを細かくできます。

\noindent 電場は、となりの点の電位の差を、距離でわって求めます（⑦の $E=-\Delta\phi/\Delta s$）。マスの大きさを\qty{1}{cm}とすると、まん中の点のたての電場は、上の点（52.7）と下の点（9.82）の差を\qty{2}{cm}でわって、$(52.7-9.82)\div0.02\approx\qty{2.1e3}{V/m}$（下向き）です。

\begin{seikaku}
マスを細かくすると、答えは本当の電位に近づきます（8-12③と同じ）。まん中の点の電位は、マスをいくつにしても、ちょうど\qty{25}{V}です。上の辺だけ\qty{100}{V}の箱と、それを90°ずつ回した3つの箱を重ね合わせると、4つの辺が全部\qty{100}{V}の箱になり、中はどこでも\qty{100}{V}になります。回しても、まん中の点は動かないので、まん中の電位は4つとも同じで、$100\div4=25$ です（8-12③のはてな？の答えと同じ考え）。
\end{seikaku}

%% ---------------------------------------------------------------
\Needspace{12\baselineskip}\bksubsection{12}{とがった所に電場が集まる}
\label{ph:電-2:sec12}
金属の表面は等電位で、電場は表面に垂直です。でも、表面の電場の\textbf{強さ}は、場所によって違います。とがった所ほど強くなるのです。

\phboxhead{2つの球をつなぐ}
半径 $R_1$ の小さな金属球と、半径 $R_2$ の大きな金属球を、細い長い導線でつなぎます。2つの球は十分離れていて、たがいの影響は無視できるとします。導線でつながっているので、2つの球は同じ電位です。
\[k\,\frac{Q_1}{R_1}=k\,\frac{Q_2}{R_2}\qquad\text{だから}\qquad\frac{Q_1}{Q_2}=\frac{R_1}{R_2}\]
表面の電場は $E=kQ/R^2$ なので、その比は
\pheq{\frac{E_1}{E_2}=\frac{Q_1/R_1^{\,2}}{Q_2/R_2^{\,2}}=\frac{R_1}{R_2}\times\frac{R_2^{\,2}}{R_1^{\,2}}=\frac{R_2}{R_1}}{単位なし}
で、\textbf{半径の小さいほうが、表面の電場が強い}のです。半径\qty{1}{cm}と\qty{10}{cm}なら、小さい球の表面の電場は10倍です。とがった先は、とても小さい半径の球のようなものなので、電場がとても強くなります。

\begin{bkfigure}
\begin{tikzpicture}[x=1mm, y=1mm]
  \fill[bkRed!15] (0,60) rectangle (120,63);
  \draw[bkRed!70!black, line width=0.6pt] (0,60) rectangle (120,63);
  \fill[bkRed!15] (39.2,30) -- (40,28.6) -- (40.8,30) -- (40.8,60) -- (39.2,60) -- cycle;
  \draw[bkRed!70!black, line width=0.6pt] (39.2,60) -- (39.2,30) -- (40,28.6) -- (40.8,30) -- (40.8,60);
  \fill[bkBlue!15] (0,-3) rectangle (120,0);
  \draw[bkBlue!70!black, line width=0.6pt] (0,-3) rectangle (120,0);
\draw[sxc lap] (118.0,6.0) -- (116.0,6.0) -- (114.0,5.9) -- (112.0,5.9) -- (110.0,5.9) -- (108.0,5.9) -- (106.0,5.9) -- (104.0,5.8) -- (102.0,5.8) -- (100.0,5.8) -- (98.0,5.8) -- (96.0,5.7) -- (94.0,5.7) -- (92.0,5.7) -- (90.0,5.6) -- (88.0,5.6) -- (86.0,5.5) -- (84.0,5.5) -- (82.0,5.4) -- (80.0,5.4) -- (78.0,5.3) -- (76.0,5.3) -- (74.0,5.2) -- (72.0,5.1) -- (70.0,5.1) -- (68.0,5.0) -- (66.0,4.9) -- (64.0,4.8) -- (62.0,4.7) -- (60.0,4.7) -- (58.0,4.6) -- (56.0,4.5) -- (54.0,4.4) -- (52.0,4.4) -- (50.0,4.3) -- (48.0,4.3) -- (46.0,4.2) -- (44.0,4.2) -- (42.0,4.2) -- (40.0,4.2) -- (38.0,4.2) -- (36.0,4.2) -- (34.0,4.3) -- (32.0,4.3) -- (30.0,4.4) -- (28.0,4.4) -- (26.0,4.5) -- (24.0,4.6) -- (22.0,4.7) -- (20.0,4.8) -- (18.0,4.9) -- (16.0,5.0) -- (14.0,5.1) -- (12.0,5.2) -- (10.0,5.4) -- (8.0,5.5) -- (6.0,5.6) -- (4.0,5.7) -- (2.0,5.9);
\draw[sxc lap] (118.0,12.0) -- (116.0,11.9) -- (114.0,11.9) -- (112.0,11.9) -- (110.0,11.8) -- (108.0,11.8) -- (106.0,11.7) -- (104.0,11.7) -- (102.0,11.6) -- (100.0,11.6) -- (98.0,11.5) -- (96.0,11.5) -- (94.0,11.4) -- (92.0,11.3) -- (90.0,11.3) -- (88.0,11.2) -- (86.0,11.1) -- (84.0,11.0) -- (82.0,10.9) -- (80.0,10.8) -- (78.0,10.7) -- (76.0,10.5) -- (74.0,10.4) -- (72.0,10.3) -- (70.0,10.1) -- (68.2,10.0) -- (68.0,10.0) -- (66.0,9.8) -- (64.0,9.7) -- (62.0,9.5) -- (60.0,9.3) -- (58.0,9.2) -- (56.0,9.0) -- (54.0,8.9) -- (52.0,8.7) -- (50.0,8.6) -- (48.0,8.5) -- (46.0,8.4) -- (44.0,8.4) -- (42.0,8.3) -- (40.0,8.3) -- (38.0,8.3) -- (36.0,8.4) -- (34.0,8.5) -- (32.0,8.6) -- (30.0,8.7) -- (28.0,8.9) -- (26.0,9.0) -- (24.0,9.2) -- (22.0,9.4) -- (20.0,9.6) -- (18.0,9.8) -- (16.4,10.0) -- (16.0,10.0) -- (14.0,10.3) -- (12.0,10.5) -- (10.0,10.7) -- (8.0,11.0) -- (6.0,11.2) -- (4.0,11.5) -- (2.0,11.7);
\draw[sxc lap] (118.0,17.9) -- (116.0,17.9) -- (114.0,17.8) -- (112.0,17.8) -- (110.0,17.7) -- (108.0,17.7) -- (106.0,17.6) -- (104.0,17.6) -- (102.0,17.5) -- (100.0,17.4) -- (98.0,17.3) -- (96.0,17.3) -- (94.0,17.2) -- (92.0,17.1) -- (90.0,17.0) -- (88.0,16.8) -- (86.0,16.7) -- (84.0,16.6) -- (82.0,16.4) -- (80.0,16.3) -- (78.0,16.1) -- (77.0,16.0) -- (76.0,15.9) -- (74.0,15.7) -- (72.0,15.5) -- (70.0,15.3) -- (68.0,15.1) -- (66.0,14.8) -- (64.0,14.6) -- (62.0,14.3) -- (60.0,14.0) -- (59.7,14.0) -- (58.0,13.8) -- (56.0,13.5) -- (54.0,13.3) -- (52.0,13.0) -- (50.0,12.8) -- (48.0,12.6) -- (46.0,12.5) -- (44.0,12.4) -- (42.0,12.3) -- (40.0,12.3) -- (38.0,12.4) -- (36.0,12.4) -- (34.0,12.6) -- (32.0,12.8) -- (30.0,13.0) -- (28.0,13.2) -- (26.0,13.5) -- (24.0,13.8) -- (22.7,14.0) -- (22.0,14.1) -- (20.0,14.4) -- (18.0,14.8) -- (16.0,15.1) -- (14.0,15.5) -- (12.0,15.8) -- (10.9,16.0) -- (10.0,16.2) -- (8.0,16.5) -- (6.0,16.9) -- (4.0,17.2) -- (2.0,17.6);
\draw[sxc lap] (118.0,23.9) -- (116.0,23.9) -- (114.0,23.8) -- (112.0,23.8) -- (110.0,23.7) -- (108.0,23.6) -- (106.0,23.5) -- (104.0,23.5) -- (102.0,23.4) -- (100.0,23.3) -- (98.0,23.2) -- (96.0,23.1) -- (94.0,23.0) -- (92.0,22.9) -- (90.0,22.7) -- (88.0,22.6) -- (86.0,22.4) -- (84.0,22.3) -- (82.0,22.1) -- (81.3,22.0) -- (80.0,21.9) -- (78.0,21.6) -- (76.0,21.4) -- (74.0,21.1) -- (72.0,20.9) -- (70.0,20.6) -- (68.0,20.2) -- (66.5,20.0) -- (66.0,19.9) -- (64.0,19.6) -- (62.0,19.2) -- (60.0,18.8) -- (58.0,18.4) -- (56.0,18.0) -- (55.8,18.0) -- (54.0,17.7) -- (52.0,17.3) -- (50.0,17.0) -- (48.0,16.7) -- (46.0,16.4) -- (44.0,16.3) -- (42.0,16.2) -- (40.0,16.1) -- (38.0,16.2) -- (36.0,16.3) -- (34.0,16.5) -- (32.0,16.8) -- (30.0,17.1) -- (28.0,17.5) -- (26.0,17.9) -- (25.7,18.0) -- (24.0,18.4) -- (22.0,18.8) -- (20.0,19.3) -- (18.0,19.8) -- (17.1,20.0) -- (16.0,20.3) -- (14.0,20.7) -- (12.0,21.2) -- (10.0,21.7) -- (8.6,22.0) -- (8.0,22.1) -- (6.0,22.6) -- (4.0,23.1) -- (2.0,23.5);
\draw[sxc lap] (118.0,29.9) -- (116.0,29.9) -- (114.0,29.8) -- (112.0,29.7) -- (110.0,29.7) -- (108.0,29.6) -- (106.0,29.5) -- (104.0,29.4) -- (102.0,29.3) -- (100.0,29.3) -- (98.0,29.1) -- (96.0,29.0) -- (94.0,28.9) -- (92.0,28.8) -- (90.0,28.6) -- (88.0,28.5) -- (86.0,28.3) -- (84.0,28.1) -- (83.2,28.0) -- (82.0,27.9) -- (80.0,27.6) -- (78.0,27.4) -- (76.0,27.1) -- (74.0,26.8) -- (72.0,26.4) -- (70.0,26.1) -- (69.7,26.0) -- (68.0,25.7) -- (66.0,25.2) -- (64.0,24.8) -- (62.0,24.3) -- (61.0,24.0) -- (60.0,23.7) -- (58.0,23.2) -- (56.0,22.6) -- (54.0,22.1) -- (53.7,22.0) -- (52.0,21.5) -- (50.0,21.0) -- (48.0,20.6) -- (46.0,20.2) -- (44.7,20.0) -- (44.0,19.9) -- (42.0,19.7) -- (40.0,19.7) -- (38.0,19.8) -- (36.0,20.0) -- (35.8,20.0) -- (34.0,20.3) -- (32.0,20.7) -- (30.0,21.2) -- (28.0,21.8) -- (27.3,22.0) -- (26.0,22.4) -- (24.0,23.1) -- (22.0,23.7) -- (21.0,24.0) -- (20.0,24.3) -- (18.0,25.0) -- (16.0,25.6) -- (14.6,26.0) -- (14.0,26.2) -- (12.0,26.8) -- (10.0,27.3) -- (8.0,27.9) -- (7.6,28.0) -- (6.0,28.4) -- (4.0,29.0) -- (2.0,29.5);
\draw[sxc lap] (118.0,35.9) -- (116.0,35.9) -- (114.0,35.8) -- (112.0,35.8) -- (110.0,35.7) -- (108.0,35.6) -- (106.0,35.5) -- (104.0,35.5) -- (102.0,35.4) -- (100.0,35.3) -- (98.0,35.2) -- (96.0,35.1) -- (94.0,34.9) -- (92.0,34.8) -- (90.0,34.6) -- (88.0,34.5) -- (86.0,34.3) -- (84.0,34.1) -- (83.2,34.0) -- (82.0,33.9) -- (80.0,33.6) -- (78.0,33.3) -- (76.0,33.0) -- (74.0,32.7) -- (72.0,32.3) -- (70.6,32.0) -- (70.0,31.9) -- (68.0,31.4) -- (66.0,30.9) -- (64.0,30.3) -- (63.1,30.0) -- (62.0,29.7) -- (60.0,29.0) -- (58.0,28.3) -- (57.3,28.0) -- (56.0,27.5) -- (54.0,26.7) -- (52.3,26.0) -- (52.0,25.9) -- (50.0,25.1) -- (48.0,24.4) -- (46.8,24.0) -- (46.0,23.8) -- (44.0,23.3) -- (42.0,22.9) -- (40.0,22.8) -- (38.0,22.9) -- (36.0,23.3) -- (34.0,23.9) -- (33.6,24.0) -- (32.0,24.6) -- (30.0,25.3) -- (28.4,26.0) -- (28.0,26.2) -- (26.0,27.1) -- (24.0,28.0) -- (23.9,28.0) -- (22.0,28.9) -- (20.0,29.7) -- (19.2,30.0) -- (18.0,30.5) -- (16.0,31.2) -- (14.0,31.9) -- (13.8,32.0) -- (12.0,32.6) -- (10.0,33.3) -- (8.0,33.8) -- (7.4,34.0) -- (6.0,34.4) -- (4.0,35.0) -- (2.0,35.5);
\draw[sxc lap] (118.0,41.9) -- (116.0,41.9) -- (114.0,41.8) -- (112.0,41.8) -- (110.0,41.7) -- (108.0,41.7) -- (106.0,41.6) -- (104.0,41.5) -- (102.0,41.5) -- (100.0,41.4) -- (98.0,41.3) -- (96.0,41.2) -- (94.0,41.1) -- (92.0,41.0) -- (90.0,40.8) -- (88.0,40.7) -- (86.0,40.5) -- (84.0,40.3) -- (82.0,40.1) -- (81.2,40.0) -- (80.0,39.9) -- (78.0,39.6) -- (76.0,39.3) -- (74.0,38.9) -- (72.0,38.6) -- (70.0,38.1) -- (69.5,38.0) -- (68.0,37.6) -- (66.0,37.1) -- (64.0,36.4) -- (62.8,36.0) -- (62.0,35.7) -- (60.0,34.9) -- (58.0,34.0) -- (58.0,34.0) -- (56.0,33.0) -- (54.2,32.0) -- (54.0,31.9) -- (52.0,30.7) -- (50.8,30.0) -- (50.0,29.5) -- (48.0,28.3) -- (47.5,28.0) -- (46.0,27.2) -- (44.0,26.3) -- (42.9,26.0) -- (42.0,25.7) -- (40.0,25.5) -- (38.0,25.8) -- (37.2,26.0) -- (36.0,26.4) -- (34.0,27.4) -- (32.8,28.0) -- (32.0,28.5) -- (30.0,29.8) -- (29.6,30.0) -- (28.0,31.1) -- (26.5,32.0) -- (26.0,32.4) -- (24.0,33.6) -- (23.2,34.0) -- (22.0,34.7) -- (20.0,35.7) -- (19.3,36.0) -- (18.0,36.6) -- (16.0,37.4) -- (14.4,38.0) -- (14.0,38.2) -- (12.0,38.9) -- (10.0,39.5) -- (8.2,40.0) -- (8.0,40.0) -- (6.0,40.6) -- (4.0,41.1) -- (2.0,41.5);
\draw[sxc lap] (118.0,48.0) -- (116.0,47.9) -- (114.0,47.9) -- (112.0,47.8) -- (110.0,47.8) -- (108.0,47.8) -- (106.0,47.7) -- (104.0,47.7) -- (102.0,47.6) -- (100.0,47.5) -- (98.0,47.5) -- (96.0,47.4) -- (94.0,47.3) -- (92.0,47.2) -- (90.0,47.1) -- (88.0,47.0) -- (86.0,46.9) -- (84.0,46.7) -- (82.0,46.6) -- (80.0,46.4) -- (78.0,46.2) -- (76.7,46.0) -- (76.0,45.9) -- (74.0,45.6) -- (72.0,45.3) -- (70.0,44.9) -- (68.0,44.5) -- (66.0,44.0) -- (66.0,44.0) -- (64.0,43.4) -- (62.0,42.7) -- (60.2,42.0) -- (60.0,41.9) -- (58.0,41.0) -- (56.3,40.0) -- (56.0,39.8) -- (54.0,38.5) -- (53.4,38.0) -- (52.0,37.0) -- (50.9,36.0) -- (50.0,35.2) -- (48.8,34.0) -- (48.0,33.3) -- (46.8,32.0) -- (46.0,31.3) -- (44.8,30.0) -- (44.0,29.4) -- (42.0,28.1) -- (41.6,28.0) -- (40.0,27.5) -- (38.4,28.0) -- (38.0,28.2) -- (36.0,29.5) -- (35.3,30.0) -- (34.0,31.4) -- (33.4,32.0) -- (32.0,33.5) -- (31.5,34.0) -- (30.0,35.6) -- (29.5,36.0) -- (28.0,37.4) -- (27.3,38.0) -- (26.0,39.1) -- (24.7,40.0) -- (24.0,40.5) -- (22.0,41.7) -- (21.3,42.0) -- (20.0,42.7) -- (18.0,43.6) -- (16.8,44.0) -- (16.0,44.3) -- (14.0,45.0) -- (12.0,45.5) -- (10.2,46.0) -- (10.0,46.1) -- (8.0,46.5) -- (6.0,46.9) -- (4.0,47.3) -- (2.0,47.7);
\draw[sxc lap] (118.0,54.0) -- (116.0,54.0) -- (114.0,53.9) -- (112.0,53.9) -- (110.0,53.9) -- (108.0,53.9) -- (106.0,53.8) -- (104.0,53.8) -- (102.0,53.8) -- (100.0,53.8) -- (98.0,53.7) -- (96.0,53.7) -- (94.0,53.6) -- (92.0,53.6) -- (90.0,53.5) -- (88.0,53.5) -- (86.0,53.4) -- (84.0,53.3) -- (82.0,53.2) -- (80.0,53.1) -- (78.0,53.0) -- (76.0,52.9) -- (74.0,52.7) -- (72.0,52.5) -- (70.0,52.3) -- (68.0,52.0) -- (67.8,52.0) -- (66.0,51.7) -- (64.0,51.4) -- (62.0,50.9) -- (60.0,50.4) -- (58.9,50.0) -- (58.0,49.7) -- (56.0,48.8) -- (54.4,48.0) -- (54.0,47.8) -- (52.0,46.4) -- (51.6,46.0) -- (50.0,44.6) -- (49.5,44.0) -- (48.0,42.1) -- (47.9,42.0) -- (46.6,40.0) -- (46.0,39.0) -- (45.5,38.0) -- (44.5,36.0) -- (44.0,35.0) -- (43.6,34.0) -- (42.7,32.0) -- (42.0,30.9) -- (41.7,30.0) -- (40.0,28.8) -- (38.3,30.0) -- (38.0,30.9) -- (37.3,32.0) -- (36.5,34.0) -- (36.0,35.2) -- (35.6,36.0) -- (34.6,38.0) -- (34.0,39.3) -- (33.6,40.0) -- (32.4,42.0) -- (32.0,42.6) -- (30.9,44.0) -- (30.0,45.1) -- (29.0,46.0) -- (28.0,46.9) -- (26.4,48.0) -- (26.0,48.3) -- (24.0,49.4) -- (22.5,50.0) -- (22.0,50.2) -- (20.0,50.9) -- (18.0,51.5) -- (16.0,51.9) -- (15.6,52.0) -- (14.0,52.3) -- (12.0,52.6) -- (10.0,52.9) -- (8.0,53.2) -- (6.0,53.4) -- (4.0,53.6) -- (2.0,53.8);
  \node[sxc lab, anchor=south] at (100,63.5) {上の電極 \qty{100}{V}};
  \node[sxc lab, anchor=north] at (100,-3.5) {下の電極 \qty{0}{V}};
  \node[sxc lab, anchor=south west, fill=white, inner sep=1pt] at (42,40) {とがった先};
  \node[sxc lab, anchor=west, fill=white, inner sep=1pt] at (121,30) {\qty{50}{V}};
\end{tikzpicture}
\bkfigcaption{とがった電極の周りの等電位線（\qty{10}{V}ごと）。ラプラス方程式を、\qty{2}{mm}ごとの点で、周りの平均にならす繰り返しで解いた。とがった先の周りで線がこみ合い、電場が強い（このマスの大きさの計算で、先のすぐ下は、右の平らな所の約5倍。マスを細かくするほど、もっと大きくなる）。}\label{ph:電-2:fig-togari}
\end{bkfigure}

\noindent 電場が、空気がこわれる強さ（約 \qty{3e6}{V/m}）をこえると、空気の分子から電子がはがれて、空気が電気を通すようになり、放電が始まります。とがった先では、ほかの所より低い電圧で、この放電が始まります。
\begin{itemize}[leftmargin=1.2\zw, itemsep=2pt, topsep=2pt, label={\color{bkTheme}●}]
\item \textbf{避雷針}：建物のいちばん高い所に、とがった金属の棒を立て、太い導線で地面とつなぎます。雷の通り道を避雷針へ導き、建物の中を電流が通らないようにします。
\item \textbf{コロナ放電}：雨の日などに、高い電圧の送電線から「ジー」という音がしたり、夜に電線の周りが青白く光ったりすることがあります。電線の表面の電場が強くなり、周りの空気がこわれかけているのです。高い電圧の送電線では、太い電線や、何本かの電線を束ねたもの、金属の輪をつけて、表面の電場を弱めています。
\item \textbf{電気集じん器}（電-1 ⑨）は、細い電線の周りのコロナ放電で、空気に電荷をつくります。
\end{itemize}

\begin{phmada}[雷は、どうやって始まるのか？]
雲の中で、氷の粒やあられがぶつかって電荷が分かれ、雲の上と下にたまることは、わかっています。ところが、雷が始まるときに雲の中で測られる電場は、大きくても1 m あたり数十万 V 程度で、乾いた空気がこわれる強さ（約 \qty{3e6}{V/m}）よりずっと小さいのです。それなのに、どうして放電が始まるのかは、まだはっきりしていません。宇宙から来る放射線（宇宙線）が空気の中でつくる高速の電子が、なだれのように数を増やして放電のきっかけを作るという説や、氷の粒のとがった所で電場が強まる（⑫）という説、その両方を組み合わせた説などが提案され、雲の中の観測と実験室の実験で調べられています。
\end{phmada}

\begin{phsaki}{コンピュータで電場を解く}
この章の表計算は、マスの数が9個だけでした。実際の機械（変圧器・送電線のがいし・モーター・電子部品）の電場を調べるときは、形を何万〜何百万個の小さな三角形や四面体に分け、ポアソン方程式を連立方程式にして解く\textbf{有限要素法}を使います。物理シリーズの『力学』（力-16）で、はりや板の変形を計算したのと同じ方法です。表面だけを分ける\textbf{境界要素法}もあります。こうした「電界解析」は、高い電圧の機械の絶縁の設計に欠かせません。
\end{phsaki}

\begin{senmonsho}[電気工学の専門書の書き方]
\begin{bkhonbun}
ガウスの法則は，電束密度 $\symbfit{D}=\varepsilon_0\symbfit{E}$ を用いて
\[ \oint_S\symbfit{D}\cdot\dd\symbfit{S}=Q,\qquad \nabla\cdot\symbfit{D}=\rho \]
と表される．電位 $V$ と電界の間には $\symbfit{E}=-\nabla V$ の関係があり，電荷のない領域で $V$ はラプラス方程式 $\nabla^2V=0$ を満たす．
\end{bkhonbun}
\bkyaku{電気工学の本は、電位を $V$、電場を「電界」とよび、$\varepsilon_0\vec E$ を\textbf{電束密度} $\vec D$〔C/m$^2$〕、その流束を\textbf{電束}〔C〕とよぶ（数学 8-8 のもっと知りたいと同じ）。両辺に $\varepsilon_0$ をかけただけで、この章の $\oiint\vec E\cdot\vec n\,\dd S=Q/\varepsilon_0$ と同じ法則。$\nabla\cdot$ は div、$\nabla$ は grad。物の中では $\vec D$ がもっと大切になる（電-3 ⑧）。}
\end{senmonsho}

\begin{taisetsu}
\def\tsl#1{\makebox[7.6\zw][l]{\sffamily\bfseries\small\color{\bkTT}#1}}
\tsl{電気力線の本数}$Q$ から $\dfrac{Q}{\varepsilon_0}=4\pi kQ$ 本。$E$＝\qty{1}{m^2}あたりの本数\\[3pt]
\tsl{ガウスの法則}$\Phi_E=\displaystyle\oiint_S\vec E\cdot\vec n\,\dd S=\frac{Q_{\text{中}}}{\varepsilon_0}$、$\mathrm{div}\,\vec E=\dfrac{\rho}{\varepsilon_0}$\\[3pt]
\tsl{電位}$\phi=k\,\dfrac{Q}{r}$、$U=q\phi$、$W=qV$（$V$ は電位の差）\\[3pt]
\tsl{一様な電場}$V=Ed$、$\qty{1}{eV}=\qty{1.60e-19}{J}$\\[3pt]
\tsl{電場と電位}$\vec E=-\mathrm{grad}\,\phi$、$\mathrm{rot}\,\vec E=\vec 0$。等電位面に垂直\\[3pt]
\tsl{ラプラス方程式}$\nabla^2\phi=-\dfrac{\rho}{\varepsilon_0}$、電荷のない所で $\nabla^2\phi=0$（周りの平均）
\end{taisetsu}

\begin{kotoba}
\small
\bkword{電場の流束}{閉じた袋から出ていく電気力線の本数。$\Phi_E$。}
\bkword{ガウスの法則}{袋から出る電気力線の本数は、中の電荷の合計の $1/\varepsilon_0$ 倍。}
\bkword{電位}{\qty{1}{C}あたりの電気の位置エネルギー。単位 V（＝J/C）。}
\bkword{電圧（電位差）}{2点の電位の差。}
\bkword{電子ボルト}{電子が\qty{1}{V}で加速されて得るエネルギー。$\qty{1.60e-19}{J}$。}
\bkword{等電位面}{電位が同じ点をつないだ面。電気力線と垂直に交わる。}
\bkword{ラプラス方程式}{電荷のない所の電位の式 $\nabla^2\phi=0$。どの点も周りの平均。}
\bkword{解の一意性}{条件が決まれば、電位の分布はただ1つに決まること。}
\end{kotoba}

\bkkigocards{\bkkigocard{$\phi$}{ファイ}{電位〔V〕。高校の教科書の電位 $V$ と同じもの。波動の本の初期位相 $\phi$ とは別の量}{electric potential}{電-2}}
{\bkkigocard{$V$}{ブイ}{電圧（2点の電位の差）〔V〕。体積 $V$ とは別の量}{voltage}{電-2}}
\bkkigocards{\bkkigocard{$\Phi_E$}{ファイ・イー}{電場の流束〔N·m$^2$/C〕。閉じた袋から出る電気力線の本数}{electric flux}{電-2}}
{\bkkigocard{$\rho$}{ロー}{電荷密度〔C/m$^3$〕。密度 $\rho$（第1・2巻）とは別の量。電-4 では $\rho_{\mathrm q}$}{charge density}{電-2}}

%% ======================== 例題 ========================
\Needspace{10\baselineskip}\bksection{reidai}{例題でたしかめよう}
\label{ph:電-2:ex1}
\begin{phreidai}[一様に帯電した球の電場]
半径\qty{0.1}{m}の球の中に、$\qty{1e-8}{C}$ の電荷が一様に広がっています。中心から\qty{0.2}{m}（球の外）と\qty{0.05}{m}（球の中）の所の電場を求めましょう。
\phbamen
対称に広がった電荷の電場を、ガウスの法則で求める場面です。
\phmushi
電荷は球の中に一様に広がっているとし、周りにほかの電荷はないとします。
\phzu
図\ref{ph:電-2:fig-taishou}(ア)。$Q=\qty{1e-8}{C}$、$R=\qty{0.1}{m}$。
\phriyu
球対称なので、半径 $r$ の球の袋の上で電場は垂直・同じ強さ。ガウスの法則で、袋の中の電荷だけを数えればよい（④）。
\phtateru
\[\text{外：}E\times4\pi r^2=\frac{Q}{\varepsilon_0},\qquad\text{中：}E\times4\pi r^2=\frac{1}{\varepsilon_0}\times Q\,\frac{r^3}{R^3}\]
\phtoku
外（$r=0.2$）：\[E=k\,\frac{Q}{r^2}=8.99\times10^9\times\frac{10^{-8}}{0.04}\approx\qty{2.2e3}{N/C}\]
中（$r=0.05$）：\[E=\frac{kQ}{R^3}\,r=\frac{89.9}{0.001}\times0.05\approx\qty{4.5e3}{N/C}\]
\phtashikame
球の中の電荷は $10^{-8}\times(0.05/0.1)^3$ で全体の $\tfrac18$ です。$\tfrac18$ の電荷が、半径0.05の所に集まっていると考えると $8.99\times10^9\times1.25\times10^{-9}\div0.0025\approx4.5\times10^3$ で、合います。\\
表面（$r=0.1$）では、どちらの式でも $\qty{8.99e3}{N/C}$ になり、つながっています（図\ref{ph:電-2:fig-kyu}）。
\kotae \qty{0.2}{m}で約 $\qty{2.2e3}{N/C}$、\qty{0.05}{m}で約 $\qty{4.5e3}{N/C}$（どちらも外向き）
\end{phreidai}
\begin{jibun}
半径\qty{0.05}{m}の球の中に $\qty{2e-8}{C}$ が一様に広がっています。中心から\qty{0.1}{m}、表面、\qty{0.025}{m}の所の電場を求めましょう。
\bkjans{\qty{0.1}{m}：約 $\qty{1.8e4}{N/C}$、表面：約 $\qty{7.2e4}{N/C}$、\qty{0.025}{m}：約 $\qty{3.6e4}{N/C}$}
\end{jibun}

\label{ph:電-2:ex2}
\begin{phreidai}[長い線の電荷を、ガウスの法則で]
1 m あたり\qty{1}{nC}の電荷が一様にある長い直線から、\qty{0.1}{m}の所の電場を、ガウスの法則で求めましょう。電-1 例題9 の答えとくらべましょう。
\phbamen
電-1 では積分で求めた電場を、袋の選び方で、楽に求める場面です。
\phmushi
直線は限りなく長く、太さを無視します。
\phzu
図\ref{ph:電-2:fig-taishou}(イ)。$\lambda=\qty{1e-9}{C/m}$、$r=\qty{0.1}{m}$、袋は半径 $r$・長さ $l$ の円柱。
\phriyu
円柱の側面では、電場は垂直・同じ強さ、上下の面では平行（④）。袋の中の電荷は $\lambda l$。
\phtateru
\[E\times2\pi rl=\frac{\lambda l}{\varepsilon_0}\qquad\text{より}\qquad E=\frac{\lambda}{2\pi\varepsilon_0r}\]
\phtoku
\[E=\frac{1\times10^{-9}}{2\pi\times8.85\times10^{-12}\times0.1}\approx\qty{180}{N/C}\]
\phtashikame
電-1 例題9 の $2k\lambda/r\approx180$ と同じ答えです。袋の長さ $l$ は、両辺で消えました。積分を1回もせずに、同じ答えが出ました。対称性がある問題では、ガウスの法則が「楽な道」です。
\kotae 約\qty{180}{N/C}（電-1 例題9 と同じ）
\end{phreidai}
\begin{jibun}
（1）1 m あたり\qty{5}{nC}の長い直線から\qty{0.3}{m}の所の電場を、ガウスの法則で求めましょう。（2）1 m$^2$ あたり\qty{2}{\micro C}の広い平面の電場を、はさむ箱の袋で求めましょう。
\bkjans{（1）約\qty{300}{N/C}。（2）$E\times2S=\sigma S/\varepsilon_0$ より、約 $\qty{1.1e5}{N/C}$}
\end{jibun}

\label{ph:電-2:ex3}
\begin{phreidai}[点電荷の周りの電位と仕事]
\qty{1}{nC}の点電荷から\qty{0.1}{m}と\qty{0.2}{m}の所の電位を求めましょう。また、\qty{2}{nC}の電荷を、\qty{0.2}{m}の所から\qty{0.1}{m}の所まで運ぶのに要る仕事を求めましょう。
\phbamen
点電荷の周りで、電荷を運ぶ仕事を電位から求める場面です。
\phmushi
点電荷とみます。運ぶ電荷は、ゆっくり動かす（運動エネルギーは変わらない）とします。
\phzu
図\ref{ph:電-2:fig-menseki}。$Q=\qty{1e-9}{C}$、$q=\qty{2e-9}{C}$、$r_{\mathrm A}=0.2$、$r_{\mathrm B}=0.1$。
\phriyu
電位は $\phi=kQ/r$（⑤）。運ぶのに要る仕事は、位置エネルギーの増え $q(\phi_{\mathrm B}-\phi_{\mathrm A})$。道によらない（保存力）。
\phtateru
\[\phi=k\,\frac{Q}{r},\qquad W=q(\phi_{\mathrm B}-\phi_{\mathrm A})\]
\phtoku
\[\phi_{\mathrm B}=\frac{8.99}{0.1}=\qty{89.9}{V},\qquad\phi_{\mathrm A}=\frac{8.99}{0.2}\approx\qty{45.0}{V}\]
\[W=2\times10^{-9}\times(89.9-45.0)\approx\qty{9.0e-8}{J}\]
\phtashikame
単位：C×V＝J。\\
大きさ：近づけるほど電位が高くなり、正の電荷どうしははね返すので、仕事が要ります（$W>0$）。反対に、\qty{0.1}{m}から\qty{0.2}{m}へ動くときは、電場が同じ大きさの仕事をして、電荷は速くなります。
\kotae 電位 \qty{89.9}{V}・約\qty{45.0}{V}。仕事 約 $\qty{9.0e-8}{J}$
\end{phreidai}
\begin{jibun}
\qty{3}{nC}の点電荷から\qty{0.2}{m}の所の電位を求めましょう。$+\qty{1}{nC}$ の電荷を、無限に遠い所からそこまで運ぶのに要る仕事はいくらですか。
\bkjans{電位 約\qty{135}{V}。仕事 約 $\qty{1.3e-7}{J}$}
\end{jibun}

\label{ph:電-2:ex4}
\begin{phreidai}[平行な板のあいだの電場]
間かく\qty{1}{cm}の2枚の平行な金属板に、\qty{100}{V}の電圧をかけます。板のあいだの電場を求めましょう。
\phbamen
一様な電場の強さを、電圧と間かくから求める場面です。
\phmushi
板は間かくに比べて十分広く、はしの電場の乱れを無視します。
\phzu
図\ref{ph:電-2:fig-ichiyou}。$V=\qty{100}{V}$、$d=\qty{0.01}{m}$。
\phriyu
一様な電場では $V=Ed$（⑥）。
\phtateru
\[E=\frac{V}{d}\]
\phtoku
\[E=\frac{100}{0.01}=\qty{1.0e4}{V/m}\]
\phtashikame
8-4⑧で、$\vec E=-\mathrm{grad}\,\phi$ から求めた値と同じです。\\
大きさ：空気がこわれる強さ（約 \qty{3e6}{V/m}）の約300分の1なので、放電は起きません。
\kotae $\qty{1.0e4}{V/m}$（N/C）、高い電位の板から低い電位の板へ向かう向き
\end{phreidai}
\begin{jibun}
間かく\qty{2}{mm}の板に、乾電池1個（\qty{1.5}{V}）をつなぎます。あいだの電場と、そこに置いた電子が受ける力の大きさを求めましょう。
\bkjans{電場 \qty{750}{V/m}、力 $1.60\times10^{-19}\times750$ で、約 $\qty{1.2e-16}{N}$}
\end{jibun}

\label{ph:電-2:ex5}
\begin{phreidai}[電子を電圧で加速する]
止まっていた電子を、\qty{100}{V}と\qty{1}{kV}の電圧で加速します。得るエネルギー（J と eV）と、速さを求めましょう（電子の質量 $\qty{9.11e-31}{kg}$）。
\phbamen
電子銃などで、電子を電圧で加速する場面です。
\phmushi
重力と、空気の分子とのぶつかりを無視します（真空の中）。相対論の効果は小さいとします。
\phzu
電荷 $e$、$m_{\mathrm e}=\qty{9.11e-31}{kg}$、$V=100$、$1000$。
\phriyu
電場がする仕事 $eV$ が、運動エネルギーになる（エネルギーの保存。⑥）。
\phtateru
\[\frac12m_{\mathrm e}v^2=eV\qquad\text{より}\qquad v=\sqrt{\frac{2eV}{m_{\mathrm e}}}\]
\phtoku
\qty{100}{V}：エネルギー $1.60\times10^{-19}\times100=\qty{1.6e-17}{J}$（\qty{100}{eV}）。
\[v=\sqrt{\frac{2\times1.6\times10^{-17}}{9.11\times10^{-31}}}\approx\qty{5.9e6}{m/s}\]
\qty{1}{kV}：エネルギー $\qty{1.6e-16}{J}$（\qty{1}{keV}）、速さは $\sqrt{10}$ 倍の約 \qty{1.9e7}{m/s}。
\phtashikame
\qty{1}{kV}の速さは光の速さ（$\qty{3.00e8}{m/s}$）の約6\,\% で、相対論の効果（原-11）は小さいと考えてよい大きさです。\\
電圧10倍で、速さは $\sqrt{10}\approx3.2$ 倍です。
\kotae \qty{100}{V}：$\qty{1.6e-17}{J}$（\qty{100}{eV}）、約 \qty{5.9e6}{m/s}。\qty{1}{kV}：$\qty{1.6e-16}{J}$（\qty{1}{keV}）、約 \qty{1.9e7}{m/s}
\end{phreidai}
\begin{jibun}
止まっていた陽子（質量 $\qty{1.67e-27}{kg}$）を\qty{1}{kV}で加速すると、エネルギーと速さはいくらですか。電子とくらべましょう。
\bkjans{エネルギーは電子と同じ $\qty{1.6e-16}{J}$（\qty{1}{keV}）。速さ 約 $\qty{4.4e5}{m/s}$（電子の約43分の1）}
\end{jibun}

\label{ph:電-2:ex6}
\begin{phreidai}[地球の表面の電荷]
晴れた日の地面の近くには、下向きに約\qty{100}{V/m}の電場があります。地面を導体の広い平面とみなして、地面 \qty{1}{m^2} あたりの電荷を求めましょう。地球全体（半径 $\qty{6.37e6}{m}$）では、何 C になりますか。
\phbamen
大気の電場から、地球の表面にある電荷を見積もる場面です。
\phmushi
地面の凹凸や、場所による電場の違いを無視し、どこでも同じ\qty{100}{V/m}とします。
\phzu
地面の上の電場 $E=\qty{100}{V/m}$（下向き：地面へ入る向き）。
\phriyu
導体の表面のすぐ外の電場は、表面の電荷 $\sigma$ だけで決まり、導体の中の電場は0なので、箱の袋の上の面だけを線が通る：$E\times S=\sigma S/\varepsilon_0$（④の平面の袋の片側だけ。電-3 ③）。
\phtateru
\[|\sigma|=\varepsilon_0E,\qquad Q=\sigma\times4\pi R^2\]
\phtoku
\[|\sigma|=8.85\times10^{-12}\times100\approx\qty{8.9e-10}{C/m^2}\]
\[|Q|=8.85\times10^{-10}\times4\pi\times(6.37\times10^6)^2\approx\qty{4.5e5}{C}\]
電場が下向き（地面に入る向き）なので、地面の電荷は\textbf{負}です。
\phtashikame
\qty{1}{m^2}あたりの電子の数は $8.9\times10^{-10}\div(1.60\times10^{-19})\approx5.5\times10^9$ 個です。地球全体の約45万 C の負の電荷とつり合う正の電荷は、空気の中と、上空の電離層にあります（コラム）。
\kotae 約 $\qty{8.9e-10}{C/m^2}$（負）。地球全体で約 $\qty{4.5e5}{C}$（負）
\end{phreidai}
\begin{jibun}
雷雲の下では、地面の近くの電場が $\qty{1e4}{V/m}$ ほどになることがあります。そのときの地面 \qty{1}{m^2} あたりの電荷の大きさを求めましょう。
\bkjans{$8.85\times10^{-12}\times10^4$ で、約 $\qty{8.9e-8}{C/m^2}$（晴れた日の100倍）}
\end{jibun}

\label{ph:電-2:ex7}
\begin{phreidai}[電位の傾きから電場を求める]
$\pm\qty{1}{nC}$ を\qty{1}{cm}離した双極子の軸の上で、まん中から\qty{0.099}{m}・\qty{0.101}{m}の所の電位を求め、⑦の $E=-\Delta\phi/\Delta s$ で\qty{0.1}{m}の所の電場を見積もりましょう。電-1 例題8 の答え（\qty{181}{N/C}）とくらべましょう。
\phbamen
電位を数値で計算し、その差から電場を求める場面です（コンピュータで電場を求めるときの基本）。
\phmushi
点電荷とみます。
\phzu
図\ref{ph:電-2:fig-toudeni}（左）の軸の上。$q=\qty{1e-9}{C}$、$d=\qty{0.01}{m}$。
\phriyu
電位はふつうのたし算で求まる（⑤）。電場は電位の下がり方（⑦）。
\phtateru
\[\phi(z)=kq\Bigl(\frac{1}{z-d/2}-\frac{1}{z+d/2}\Bigr),\qquad E\approx-\frac{\phi(0.101)-\phi(0.099)}{0.002}\]
\phtoku
\[\phi(0.099)=8.99\Bigl(\frac{1}{0.094}-\frac{1}{0.104}\Bigr)\approx\qty{9.196}{V},\qquad\phi(0.101)\approx\qty{8.835}{V}\]
\[E\approx-\frac{8.835-9.196}{0.002}\approx\qty{181}{N/C}\]
\phtashikame
電-1 例題8 の重ね合わせの答え 181 と、ぴったり合います。電位は「向きのない数」なので、矢印のたし算より計算が楽です。\\
電位は\qty{0.1}{m}で約\qty{9.0}{V}で、点電荷の\qty{89.9}{V}の約10分の1です（双極子の電位は $1/r^2$ で弱まる）。
\kotae 電位 約\qty{9.20}{V}・約\qty{8.83}{V}。電場 約\qty{181}{N/C}（電-1 例題8 と同じ）
\end{phreidai}
\begin{jibun}
\qty{1}{nC}の点電荷から\qty{0.099}{m}・\qty{0.101}{m}の所の電位を求め、\qty{0.1}{m}の所の電場を見積もりましょう。$E=kQ/r^2$ とくらべましょう。
\bkjans{電位 約\qty{90.81}{V}・約\qty{89.01}{V}。電場 $(90.81-89.01)\div0.002$ で、約\qty{899}{N/C}（$kQ/r^2=899$ と同じ）}
\end{jibun}

\label{ph:電-2:ex8}
\begin{phreidai}[電位が0になる点]
$+\qty{2}{nC}$ の点電荷を $x=0$ に、$-\qty{1}{nC}$ の点電荷を $x=\qty{0.3}{m}$ に置きます。$x$ 軸の上で、電位が0になる点を全部求めましょう。
\phbamen
2つの電荷の電位の和が0になる点を探す、高校の入試にもよく出る形です。
\phmushi
点電荷とみます。無限に遠い所の電位を0とします。
\phzu
求める点を $x$ とし、$+2$\,nC までの距離 $|x|$、$-1$\,nC までの距離 $|x-0.3|$。
\phriyu
電位は向きのない数なので、2つの電位をふつうにたす（⑤）。正の電荷の電位は正、負の電荷の電位は負。
\phtateru
\[k\,\frac{2\times10^{-9}}{|x|}-k\,\frac{1\times10^{-9}}{|x-0.3|}=0\qquad\text{つまり}\qquad\frac{2}{|x|}=\frac{1}{|x-0.3|}\]
\phtoku
あいだ（$0<x<0.3$）：$2(0.3-x)=x$ より $x=\qty{0.2}{m}$。\\
$-1$\,nC の外（$x>0.3$）：$2(x-0.3)=x$ より $x=\qty{0.6}{m}$。\\
$+2$\,nC の外（$x<0$）：$2|x-0.3|=|x|$ は、$x<0$ では $2(0.3-x)=-x$ から $x=0.6$ となって $x<0$ に合わないので、答えはありません。
\phtashikame
$x=0.2$：$2/0.2=10$、$1/0.1=10$ で等しい。$x=0.6$：$2/0.6\approx3.33$、$1/0.3\approx3.33$ で等しい。\\
\textbf{電位が0の点では、電場は0ではありません}。たとえば $x=0.2$ では、どちらの電荷の電場も $+x$ の向きで、たし合わされています。電位0と電場0は別のことです（まちがえやすいところ）。
\kotae $x=\qty{0.2}{m}$ と $x=\qty{0.6}{m}$
\end{phreidai}
\begin{jibun}
$+\qty{4}{nC}$ を $x=0$ に、$-\qty{1}{nC}$ を $x=\qty{0.5}{m}$ に置きます。$x$ 軸の上で電位が0になる点を全部求めましょう。
\bkjans{あいだ：$4(0.5-x)=x$ より $x=\qty{0.4}{m}$。外：$4(x-0.5)=x$ より $x=\tfrac23\approx\qty{0.67}{m}$}
\end{jibun}

\label{ph:電-2:ex9}
\begin{phreidai}[小さな金属球は、どこまで電圧をかけられるか]
半径\qty{1}{cm}の金属球を、\qty{30}{kV}の電位にしました。表面の電場を求め、空気がこわれる強さ（約 $\qty{3e6}{V/m}$）とくらべましょう。
\phbamen
とがった所・小さな球で放電が起きやすいわけを、数で確かめる場面です。
\phmushi
球の周りにほかの物はなく、球の電荷は表面に一様に広がっているとします。
\phzu
半径 $R=\qty{0.01}{m}$、電位 $\phi=\qty{3.0e4}{V}$。
\phriyu
球の外の電場と電位は、点電荷と同じ（④⑤）。$\phi=kQ/R$、$E=kQ/R^2$ なので、$E=\phi/R$。
\phtateru
\[E=\frac{kQ}{R^2}=\frac{\phi}{R}\]
\phtoku
\[E=\frac{3.0\times10^4}{0.01}=\qty{3.0e6}{V/m}\]
電荷は $Q=\phi R/k=3.0\times10^4\times0.01\div(8.99\times10^9)\approx\qty{3.3e-8}{C}$ です。
\phtashikame
ちょうど空気がこわれる強さになり、放電（パチッという火花やコロナ放電）が起き始めます。同じ\qty{30}{kV}でも、半径\qty{10}{cm}の球なら表面は $3\times10^5$\,V/m で、放電しません。小さい球・とがった先ほど、低い電圧で放電します（⑫）。
\kotae 約 $\qty{3.0e6}{V/m}$。空気がこわれる強さと同じくらいで、放電が始まる
\end{phreidai}
\begin{jibun}
空気がこわれる強さを $\qty{3e6}{V/m}$ として、半径\qty{5}{mm}の金属球と、半径\qty{10}{cm}の金属球が、放電せずにいられる電位の上限を求めましょう。
\bkjans{\qty{5}{mm}：$3\times10^6\times0.005=\qty{1.5e4}{V}$（\qty{15}{kV}）。\qty{10}{cm}：\qty{3.0e5}{V}（\qty{300}{kV}）}
\end{jibun}

\label{ph:電-2:ex10}
\begin{phreidai}[同軸ケーブルの中の電場]
テレビのアンテナ線のような\textbf{同軸ケーブル}は、半径 $a$ の中心の導線と、それを囲む半径 $b$ の円筒の導体からできています。$a=\qty{0.5}{mm}$、$b=\qty{2}{mm}$ で、あいだが空気、中心の導線と外の円筒のあいだに\qty{1}{kV}の電圧をかけるとき、中心の導線の表面の電場を求めましょう。
\phbamen
円柱対称な電極のあいだの電場を、ガウスの法則と電位から求める場面です。
\phmushi
ケーブルは十分長く、はしの効果を無視します。あいだは真空と同じとします。
\phzu
中心の導線に1 m あたり $\lambda$ の電荷。$a=\qty{5e-4}{m}$、$b=\qty{2e-3}{m}$、$V=\qty{1000}{V}$。
\phriyu
円柱の袋で $E=\lambda/(2\pi\varepsilon_0r)$（④）。電圧は、電場を $a$ から $b$ まで積分したもの（⑤⑨）。2つの式から $\lambda$ を消す。
\phtateru
\[V=\int_a^b\frac{\lambda}{2\pi\varepsilon_0r}\,\dd r=\frac{\lambda}{2\pi\varepsilon_0}\ln\frac{b}{a},\qquad E(a)=\frac{\lambda}{2\pi\varepsilon_0a}=\frac{V}{a\ln(b/a)}\]
\phtoku
\[\ln\frac{b}{a}=\ln4\approx1.386,\qquad E(a)=\frac{1000}{5\times10^{-4}\times1.386}\approx\qty{1.4e6}{V/m}\]
\phtashikame
外の円筒の内側では、半径が4倍なので $\tfrac14$ の約 \qty{3.6e5}{V/m} です。中心の細い導線の周りに電場が集まります（⑫）。空気がこわれる強さ（約 \qty{3e6}{V/m}）の半分くらいなので、もっと電圧を上げると、中心の導線の周りから放電が始まります。本物のケーブルは、あいだにポリエチレンなどの絶縁体をつめて、こわれにくくしています（電-3）。
\kotae 約 $\qty{1.4e6}{V/m}$
\end{phreidai}
\begin{jibun}
$a=\qty{1}{mm}$、$b=\qty{4}{mm}$ の同軸ケーブルに\qty{1}{kV}をかけるとき、中心の導線の表面と、外の円筒の内側の電場を求めましょう。
\bkjans{$\ln4\approx1.386$。中心の表面：約 $\qty{7.2e5}{V/m}$、外の内側：約 $\qty{1.8e5}{V/m}$}
\end{jibun}

\label{ph:電-2:ex11}
\begin{phreidai}[表計算で箱の中の電位を求める]
⑪の正方形の箱（上の辺\qty{100}{V}、ほかの辺\qty{0}{V}、$5\times5$ の点）で、ならす繰り返しの1回めと2回めの、上の段とまん中の段の電位を求めましょう。また、落ち着いたときに、まん中の点が「周り4点の平均」になっていることを確かめましょう。
\phbamen
ラプラス方程式を、表計算の繰り返しで解く場面です。
\phmushi
マスの大きさより細かい電位の変化を無視します。箱の中に電荷はありません。
\phzu
表\ref{ph:電-2:tab-sheet}。はじめは、まん中の9点が\qty{0}{V}。
\phriyu
電荷のない所では、電位は周りの平均（⑩）。繰り返すと、本当の答えに近づく（⑪）。
\phtateru
\[\phi_{\text{新}}=\frac{\phi_{\text{北}}+\phi_{\text{西}}+\phi_{\text{東}}+\phi_{\text{南}}}{4}\]
\phtoku
1回め：上の段は $(100+0+0+0)\div4=25$（3点とも）、まん中の段は0。\\
2回め：上の段の左は $(100+0+25+0)\div4=31.25$、まん中は $(100+25+25+0)\div4=37.5$、右は31.25。まん中の段は $(25+0+0+0)\div4=6.25$。\\
落ち着いたとき（表\ref{ph:電-2:tab-kotae}）：まん中の点の周りは、北52.679、西18.75、東18.75、南9.821 なので、
\[\frac{52.679+18.75+18.75+9.821}{4}=\frac{100}{4}=25\]
\phtashikame
繰り返すほど、上のふちの\qty{100}{V}が、少しずつ下へしみこんでいくようすがわかります。8-12③の板の温度と、まったく同じ数です。電位も温度も、同じラプラス方程式をみたすからです。
\kotae 1回め：上の段 25・25・25、まん中の段 0。2回め：上の段 31.25・37.5・31.25、まん中の段 6.25。落ち着いたまん中は\qty{25}{V}
\end{phreidai}
\begin{jibun}
同じ $5\times5$ の箱で、（1）上の辺と左の辺を\qty{100}{V}、ほかを\qty{0}{V}にしたとき、（2）上の辺と下の辺を\qty{100}{V}、ほかを\qty{0}{V}にしたときの、まん中の点の電位を、⑪のせいかくに言うとの「重ね合わせ」で求めましょう。
\bkjans{（1）上だけ\qty{25}{V}と左だけ\qty{25}{V}の重ね合わせで\qty{50}{V}。（2）同じく\qty{50}{V}}
\end{jibun}

\label{ph:電-2:ex12}
\begin{phreidai}[閉じた金属の箱の中]
⑪の箱で、4つの辺を全部\qty{50}{V}にしました。箱の中の電位と電場を、表計算と、解の一意性の2つの方法で求めましょう。
\phbamen
周りを金属で囲まれた所の電場を考える場面です（電-3 の静電遮へいの入口）。
\phmushi
箱の中に電荷はないとします。
\phzu
$5\times5$ の点、ふちは全部\qty{50}{V}。
\phriyu
表計算では周りの平均（⑪）。式では、ラプラス方程式の答えは1つ（⑩）。
\phtateru
表計算：$\phi_{\text{新}}=$（北＋西＋東＋南）÷4。一意性：「どこでも $\phi=50$」が答えになるかを確かめる。
\phtoku
表計算：はじめ0から始めると、1回めの上の段の左は $(50+50+0+0)\div4=25$ などと増えていき、繰り返すと、まん中の9点は全部\qty{50}{V}に落ち着きます。\\
一意性：「どこでも\qty{50}{V}」は、どの点も周りの平均（50）と同じで、ふちも\qty{50}{V}なので、条件を全部みたします。答えは1つなので、これが答えです。電位が一定なので、電場は0です。
\phtashikame
はじめの値を0にしても100にしても、繰り返すと同じ\qty{50}{V}に落ち着きます（答えは1つ）。箱が何 V であっても、中の電場は0です。
\kotae 箱の中はどこでも\qty{50}{V}、電場は0
\end{phreidai}
\begin{jibun}
同じ箱の4つの辺を全部\qty{120}{V}にすると、中の電位と電場はどうなりますか。まん中の点の電位を、表計算の考えで確かめましょう。
\bkjans{どこでも\qty{120}{V}、電場は0。まん中の点も $(120+120+120+120)\div4=120$}
\end{jibun}

%% ======================== まちがえやすいところ ========================
\Needspace{11\baselineskip}\bksection{tsumazuki}{まちがえやすいところをチェック}
\begin{tsumazuki}{電位が0なら、電場も0？}
\bkng{例題8の $x=\qty{0.2}{m}$ は電位が0なので、そこでは電荷は力を受けない。}
\bkok{電位0は「高さ0」、電場は「坂の急さ」。高さ0の所でも、坂が急なら力を受ける。}
電場は電位の\textbf{下がり方}で決まります（⑦）。電位の値そのものは、基準の決め方で変わります。
\end{tsumazuki}

\begin{tsumazuki}{電場が0なら、電位も0？}
\bkng{同じ大きさの正の電荷2つのまん中は、電場が0なので、電位も0。}
\bkok{まん中の電位は、2つの電荷の電位の和 $2kQ/r$ で、0ではない。電場が0なのは、そこで電位の傾きが0（山の尾根の峠のような所）だから。}
図\ref{ph:電-2:fig-toudeni}（右）のまん中の点は、電位が高い所の中で、傾きが0になっている所です。電場は電位の\textbf{傾き}、電位は\textbf{高さ}です（⑦）。
\end{tsumazuki}

\begin{tsumazuki}{電位も矢印でたす？}
\bkng{2つの電荷の電位を、向きを考えて成分に分けてたす。}
\bkok{電位は向きのない数なので、そのままたす（正の電荷は正、負の電荷は負）。}
矢印でたすのは電場です。電位でたしてから傾きで電場を出すと、計算が楽になります（例題7）。
\end{tsumazuki}

\begin{tsumazuki}{負の電荷も、電位の低いほうへ動く？}
\bkng{電荷は、電位の高い所から低い所へ動く。}
\bkok{正の電荷は電位の低いほうへ、負の電荷（電子）は電位の\textbf{高い}ほうへ動こうとする。}
位置エネルギー $U=q\phi$ が下がる向きに動きます。$q<0$ なら、$\phi$ が高いほど $U$ は低くなります（⑤）。
\end{tsumazuki}

\begin{tsumazuki}{ガウスの法則で、いつでも電場が出せる？}
\bkng{どんな電荷でも、袋を選べば、ガウスの法則だけで電場が求まる。}
\bkok{法則はいつも正しいが、「電場が面に垂直で同じ強さ」の袋が見つかるのは、対称なときだけ。}
対称でないときは、重ね合わせの積分（電-1 ⑦）か、ラプラス方程式を数値で解きます（⑪）。
\end{tsumazuki}

\begin{tsumazuki}{電子ボルトは電圧？}
\bkng{1\,eV は1\,V と同じ大きさの電圧。}
\bkok{eV はエネルギーの単位。電荷 $e$ の粒が\qty{1}{V}で加速されて得るエネルギー $\qty{1.60e-19}{J}$。}
「ボルト」の名前がついていますが、J と同じ種類の量です（⑥）。
\end{tsumazuki}

\begin{column}{地球は、大きな電気の入れ物}
例題6で、晴れた日の地面には、全体で約45万 C の負の電荷があることがわかりました。上空およそ60 km より上には、太陽の光や宇宙線で空気がイオンになった\textbf{電離層}があり、電気をよく通します。地面と電離層は、空気をはさんだ2枚の板のようなもので、電離層は地面より約30万 V 高い電位になっていると見積もられています。空気は少しだけ電気を通すので、晴れた所では、電離層から地面へ、とても弱い電流が流れ続けています。それでも電荷がなくならないのは、世界のどこかでいつも起きている雷雨が、地面へ負の電荷を運び、電荷をくみ上げ続けているからだと考えられています。地球全体を1つの回路とみるこの考えを「全地球の電気回路」とよび、雷や大気の研究が続けられています。
\end{column}

%% ======================== 工学ではこう使う ========================
\Needspace{18\baselineskip}\bksection{phgear}{工学ではこう使う}
\begin{phkougaku}[電界解析と絶縁設計]
\textbf{高い電圧の機械の形}：変電所の機械や送電線の金具は、角を丸くし、表面をなめらかにしてあります。とがった所や細い所は、電場が集まって（⑫）放電が起きやすくなるからです。設計の段階で、コンピュータでラプラス方程式を解いて（⑪の方法を細かくしたもの）、表面の電場がいちばん強くなる所を探し、空気や絶縁物がこわれる強さより十分小さくなるように、形を決めます。

\textbf{同軸ケーブル}：例題10のように、中心の導線の周りに電場が集まります。外の円筒の半径を決めたとき、中心の導線を細くしすぎても太くしすぎても、表面の電場は強くなります。いちばん弱くなるのは $\ln(b/a)=1$ のときです（章末問題）。

\textbf{電子を電圧で動かす機械}：電子顕微鏡、X 線の装置、昔のブラウン管は、電圧で電子を加速し（⑥）、電場で曲げて（電-1 ⑧）、行き先を決めます。

\begin{phketa}[電場と電位の数の目安]
\phketaline{晴れた日の地面の近くの電場}{約\qty{100}{V/m}（下向き）}
\phketaline{空気がこわれる電場（1気圧）}{約 \qty{3e6}{V/m}（\qty{3}{kV/mm}）}
\phketaline{乾電池・家のコンセント}{\qty{1.5}{V}・\qty{100}{V}（実効値：電-10）}
\phketaline{\qty{1}{nC}から\qty{0.1}{m}の電位}{\qty{89.9}{V}}
\phboxhead{答えが桁で合っているか}
電場が $10^6$\,V/m をこえる答えが出たら、空気の中なら放電が起きていないか考えましょう。電位と電場の単位の関係（V ＝ V/m × m）で、距離を cm や mm のまま入れていないかも確かめます。
\end{phketa}
\end{phkougaku}
\phstation{電気磁気学,高電圧工学,電界解析,電子デバイス}{電位、電位差、等電位面、ガウスの法則、電束、電界集中、コロナ放電、ラプラス方程式、有限要素法}

%% ======================== 章末問題 ========================
\bkendproblems
\bkproblemgroup{基本}{3}
\bkproblem 半径\qty{0.1}{m}の金属球に\qty{5}{nC}の電荷をあたえました（電荷は表面に一様に広がる）。（1）中心から\qty{0.3}{m}の所の電場と電位を求めましょう。（2）球の表面の電位を求めましょう。
\bkpans{（1）$E=8.99\times5\div0.09$ で、約\qty{500}{N/C}。電位 約\qty{150}{V}。（2）約\qty{450}{V}}
\bkproblem 間かく\qty{5}{mm}の平行な板に\qty{12}{V}をかけました。（1）あいだの電場を求めましょう。（2）$-$ の板の近くで止まっていた電子が $+$ の板まで加速されると、得るエネルギー（J と eV）と速さは？
\bkpans{（1）\qty{2400}{V/m}。（2）\qty{12}{eV}＝ $\qty{1.9e-18}{J}$、速さ約 $\qty{2.1e6}{m/s}$}
\bkproblem 電位\qty{300}{V}の点 A から、電位\qty{100}{V}の点 B まで、$+\qty{2}{\micro C}$ の電荷が動きました。（1）電場の力がした仕事を求めましょう。（2）道すじを変えると、仕事は変わりますか。わけも答えましょう。
\bkpans{（1）$2\times10^{-6}\times(300-100)=\qty{4.0e-4}{J}$。（2）変わらない。電場の力は保存力で、仕事は電位の差だけで決まるから}
\bkproblem $+\qty{3}{nC}$ を $x=0$ に、$-\qty{1}{nC}$ を $x=\qty{0.4}{m}$ に置きます。$x$ 軸の上で電位が0になる点を全部求めましょう。
\bkpans{あいだ：$3(0.4-x)=x$ より $x=\qty{0.3}{m}$。外：$3(x-0.4)=x$ より $x=\qty{0.6}{m}$}
\bkproblem 晴れた日の地面の近くの電場を\qty{100}{V/m}とすると、身長\qty{1.7}{m}の人の頭の高さと足もとの電位の差は、何もないときなら何 V になりますか。それでも感電しないのはなぜか、⑦の考えで説明しましょう。
\bkpans{$100\times1.7=\qty{170}{V}$。人の体は電気を通すので、立つと体全体がほぼ1つの等電位（地面と同じ）になり、周りの等電位面が体をよけるように曲がる。体の中に大きな電圧はかからない}
\bkproblem \qty{4}{nC}の点電荷から\qty{0.3}{m}の所の電位を求めましょう。そこに $+\qty{2}{nC}$ の電荷を置いたときの位置エネルギーはいくらですか（無限に遠い所を0とする）。
\bkpans{電位 $8.99\times4\div0.3$ で、約\qty{120}{V}。位置エネルギー $2\times10^{-9}\times120$ で、約 $\qty{2.4e-7}{J}$}
\bkproblemgroup{チャレンジ}{5}
\bkproblem 一様に帯電した球（電荷 $Q$、半径 $R$）の中の電場 $E=kQr/R^3$ から、⑧の発散の式で、球の中の電荷の密度 $\rho$ を求めましょう。$Q=\qty{1e-8}{C}$、$R=\qty{0.1}{m}$ のとき、$\rho$ はいくらですか。
\bkpans{$\mathrm{div}\,\vec E=3kQ/R^3=\rho/\varepsilon_0$ より $\rho=3Q/(4\pi R^3)$（電荷÷球の体積）。約 $\qty{2.4e-6}{C/m^3}$}
\bkproblem ⑪の箱を $7\times7$ の点に細かくして、Python でならす繰り返しを行い、まん中の点の電位を求めましょう。その値になるわけも説明しましょう。
\bkpans{\qty{25}{V}。上・左・下・右の辺を\qty{100}{V}にした4つの箱を重ね合わせると、中はどこでも\qty{100}{V}。4つの箱は回したものどうしで、まん中の電位は同じなので、$100\div4=25$（マスの数によらない）}
\bkproblem 半径\qty{1}{cm}と\qty{10}{cm}の金属球を、細い長い導線でつなぎ、全部で $\qty{1.1e-7}{C}$ の電荷をあたえました。2つの球は十分離れているとして、それぞれの電荷、電位、表面の電場を求めましょう。
\bkpans{電位が同じなので電荷は半径に比例し、小さい球 $\qty{1e-8}{C}$、大きい球 $\qty{1e-7}{C}$。電位はどちらも約 $\qty{9.0e3}{V}$。表面の電場は、小さい球 約 $\qty{9.0e5}{V/m}$、大きい球 約 $\qty{9.0e4}{V/m}$（10倍の違い）}
\bkproblem 外の円筒の半径 $b=\qty{2}{mm}$ の同軸ケーブルに\qty{1}{kV}をかけます。中心の導線の表面の電場（例題10）をいちばん小さくする半径 $a$ を、$a\ln(b/a)$ の微分（5-7）で求めましょう。そのときの電場はいくらですか。
\bkpans{$\dfrac{\dd}{\dd a}\bigl(a\ln b-a\ln a\bigr)=\ln\dfrac{b}{a}-1=0$ より $a=b/e\approx\qty{0.74}{mm}$。電場は $V/a$ で、約 $\qty{1.4e6}{V/m}$}

\begin{furikaeri}
\begin{checklist}
\item ガウスの法則で、球・長い直線・広い平面の電荷の電場を求められる。\bkfurref{ph:電-2:sec4}
\item 電位を使って、電荷を運ぶ仕事や、加速された電子の速さを求められる。\bkfurref{ph:電-2:sec5}
\item 等電位面と電気力線の関係を読み、一様な電場で $V=Ed$ を使える。\bkfurref{ph:電-2:sec7}
\item ガウスの法則の微分形と $\vec E=-\mathrm{grad}\,\phi$ を、ベクトル解析で説明できる。\bkfurref{ph:電-2:sec8}
\item ラプラス方程式を表計算で解き、電場が集まる所を示せる。\bkfurref{ph:電-2:sec11}
\end{checklist}
\end{furikaeri}
\bkfukidashi{sora}{電気の力が距離の2乗で弱まるのは、空間が3次元だからなんだ。光や音と同じ理由だったなんて、びっくり！}
